%%%%%%%%%%%%%%%%%%%%%%%%%%%%%%%%%%%%%%%%%%%%%%%%%%%%%%%%%%%%%%%%%%%%%%
% Lớp 11 | Giới hạn. Hàm số liên tục | Giới hạn của dãy số
% Dạng 1D3?1-3: Phương pháp lượng liên hợp (lim hữu hạn)
% Số câu: 95  (sinh tự động bởi tach_cau.py; từ nay thầy sửa trực tiếp file này)
%%%%%%%%%%%%%%%%%%%%%%%%%%%%%%%%%%%%%%%%%%%%%%%%%%%%%%%%%%%%%%%%%%%%%%

%%%%%%%%%%%%%%%% TRẮC NGHIỆM 4 PHƯƠNG ÁN %%%%%%%%%%%%%%%%

% nguồn: Ok_Toan11_TracNghiem_Lan2
\begin{ex}%[1D3N1-3]%[ma:e9d5987a]
	Tính  giới hạn $L=\lim\left(\sqrt{4n^2+n}-\sqrt{4n^2+2} \right)$.
	 \choice
		 {$+\infty$}
		{$-7$} 
		{$-\infty$} 
		 {\True $\dfrac{1}{4}$}
	\loigiai{ 
		\begin{align*}
			L&=\lim\left(\sqrt{4n^2+n}-\sqrt{4n^2+2} \right)\\ 
			&=\lim\dfrac{\left(\sqrt{4 n^2+n}-\sqrt{4  n^2+2}\right)\left(\sqrt{4 n^2+n}+\sqrt{4 n^2+2}\right)}{\sqrt{4 n^2+n}+\sqrt{4 n^2+2}}\\
			&=\lim\dfrac{\left(4 n^2+n\right)-\left(4 n^2+2\right)}{\sqrt{4 n^2+n}+\sqrt{4  n^2+2}}\\
			 &=\lim\dfrac{n-2}{\sqrt{4 n^2+n}+\sqrt{4 n^2+2}}\\
			&=\lim\dfrac{n\left(1-\dfrac{2}{n}\right)}{n\left(\sqrt{4+\dfrac{1}{n}}+\sqrt{4+\dfrac{2}{n^2}}\right)}\\ 
			&=\lim\dfrac{1-\dfrac{2}{n}}{\sqrt{4+\dfrac{1}{n}}+\sqrt{4+\dfrac{2}{n^2}}}\\ 
			&=\dfrac{1-0}{\sqrt{4+0}+\sqrt{4+0}}\\ 
			&=\dfrac{1}{4}. 
		\end{align*}
	} 
\end{ex}

% nguồn: Ok_Toan11_TracNghiem_Lan2
\begin{ex}%[1D3N1-3]%[ma:d0c24cf3]
	Tính  giới hạn $ L=\lim\left(\sqrt[3]{8n^3+3n^2-2}+ \sqrt[3]{5n^2- 8n^3} \right)$.
	\choice 
		{$+\infty$} 
		{$-7$} 
		{$\dfrac{53}{2}$} 
		{\True  $\dfrac{2}{3}$}
	 \loigiai{
		\begin{align*}
			L& =\lim\left(\sqrt[3]{8  n^3+3 n^2-2}+\sqrt[3]{5 n^2-8 n^3}\right) \\
			& =\lim\dfrac{8  n^2-2}{\sqrt[3]{\left(8 n^3+3 n^2-2\right)^2}-\sqrt[3]{\left(8 n^3+3 n^2-2\right) \left(5 n^2-8 n^3\right)}+\sqrt[3]{\left(5 n^2-8 n^3\right)^2}} \\
			& =\lim\dfrac{8-\dfrac{2}{8 n^2}}{\sqrt[3]{\left(8+\dfrac{3}{n}-\dfrac{2}{n^3}\right)^2}-\sqrt[3]{\left(8+\dfrac{3}{n}-\dfrac{2}{n^3}\right) \left(\dfrac{5}{n}-8\right)}+\sqrt[3]{\left(\dfrac{5}{n}-8\right)^2}}\\ 
			&=\dfrac{2}{3} . 
		\end{align*}
	} 
\end{ex}

% nguồn: Ok_Toan11_TracNghiem_Lan2
\begin{ex}%[1D3N1-3]%[ma:09aea189]
	 Tính giới hạn $\lim\left(n-\sqrt{n^2-4n} \right)$.
	 \choice
		{$3$} 
		 {$1$}
		{\True $2$} 
		{$4$} 
	\loigiai{ 
		 Ta có
		\begin{align*}
			\lim\left(n-\sqrt{n^2-4n} \right) 
			 &=\lim\dfrac{\left(n-\sqrt{n^2-4n} \right)\left(n+\sqrt{n^2-4n} \right)}{n+\sqrt{n^2-4n}}\\
			&=\lim\dfrac{4n}{n+\sqrt{n^2-4n}}\\ 
			&=\lim\dfrac{4}{1+\sqrt{1-\dfrac{4}{n}}}\\ 
			 &=2.
		\end{align*}
	} 
\end{ex}

% nguồn: Ok_Toan11_TracNghiem_Lan2
\begin{ex}%[1D3N1-3]%[ma:6b7ed02b]
	Giới hạn $\lim\sqrt{n}  \left(\sqrt{n+4}-\sqrt{n+3} \right)$ bằng
	 \choice
		{$0$} 
		 {$+\infty$}
		{$\dfrac{7}{2}$} 
		 {\True $\dfrac{1}{2}$}
	\loigiai{ 
		\begin{align*}
			 \lim\sqrt{n} \left(\sqrt{n+4}-\sqrt{n+3} \right)
			&=\lim\sqrt{n} \dfrac{\left(\sqrt{n+4}-\sqrt{n+3} \right) \left(\sqrt{n+4}+\sqrt{n+3}  \right)}{\sqrt{n+4}+\sqrt{n+3}}\\
			&=\lim\sqrt{n} \dfrac{1}{\sqrt{n+4}+\sqrt{n+3}}\\ 
			 &=\lim\dfrac{1}{\sqrt{1+\dfrac{4}{n}}+\sqrt{1+\dfrac{3}{n}}}\\
			 &=\dfrac{1}{2}.
		\end{align*}
	} 
\end{ex}

% nguồn: Ok_Toan11_TracNghiem_Lan2
\begin{ex}%[1D3N1-3]%[ma:c99b60c7]
	Trong các giới hạn sau đây,  giới hạn nào có giá trị bằng $1$?
	 \choice
		 {$\lim\dfrac{3^{n+1}+2n}{5+3^n}$}
		{$\lim\dfrac{3n^2+n}{4n^2-5}$} 
		{\True  $\lim\left(\sqrt{n^2+2n}-\sqrt{n^2+1} \right)$}
		 {$\lim\dfrac{2n^3+3}{1+2n^2}$}
	 \loigiai{
		Ta  có
		\begin{align*}
			\lim\left(\sqrt{n^2+2n}-\sqrt{n^2+1} \right) 
			 &=\lim\dfrac{\left(\sqrt{n^2+2n}-\sqrt{n^2+1} \right)\left(\sqrt{n^2+2n}+\sqrt{n^2+1} \right)}{\sqrt{n^2+2n}+\sqrt{n^2+1}}\\
			&=\lim\dfrac{2n-1}{\sqrt{n^2+2n}+\sqrt{n^2+1}}\\ 
			&=\lim\dfrac{2-\dfrac{1}{n}}{\sqrt{\dfrac{n^2+2n}{n^2}}+\sqrt{\dfrac{n^2+1}{n^2}}}\\ 
			&=\lim\dfrac{2-\dfrac{1}{n}}{\sqrt{1+\dfrac{2}{n}}+\sqrt{1+\dfrac{1}{n}}}\\ 
			&=1. 
		\end{align*}
	 }
\end{ex}

% nguồn: Ok_Toan11_TracNghiem_Lan2
\begin{ex}%[1D3N1-3]%[ma:ef7a951e]
	 Giới hạn $\lim\left(\sqrt{n^2+18n}-n\right)$ bằng
	\choice 
		 {\True $9$}
		 {$+\infty$}
		 {$18$}
		 {$0$}
	\loigiai{ 
		\begin{align*}
			 \lim\left(\sqrt{n^2+18n}-n\right)
			 &=\lim\dfrac{\left(\sqrt{n^2+18n}-n\right)\left(\sqrt{n^2+18n}+n\right)}{\sqrt{n^2+18n}+n}\\
			 &=\lim\dfrac{18n}{\sqrt{n^2+18n}+n}\\
			&=\lim\dfrac{18}{\sqrt{1+\dfrac{18}{n}}+1}\\ 
			&=9. 
		\end{align*}
	 }
\end{ex}

% nguồn: Ok_Toan11_TracNghiem_Lan2
\begin{ex}%[1D3N1-3]%[ma:0282ea4f]
	Tính $\lim\limits_{x\to 2} \dfrac{\sqrt{x+2}-2}{x-2}$. 
	\choice 
		{$\dfrac{1}{2}$} 
		{\True  $\dfrac{1}{4}$}
		 {$0$}
		 {$1$}
	 \loigiai{
		 $\lim\limits_{x\to 2} \dfrac{\sqrt{x+2}-2}{x-2}
		=  \lim\limits_{x\to 2} \dfrac{x-2}{(x-2)(\sqrt{x+2}+2)}
		= \lim\limits_{x\to 2} \dfrac{1}{\sqrt{x+2}+2} 
		=  \dfrac{1}{4}$.
	 }
\end{ex}

% nguồn: Ok_Toan11_TracNghiem_Lan2
\begin{ex}%[1D3N1-3]%[ma:016f2221]
	Tính giới hạn sau $L=\lim\left(\sqrt[3]{n+4}  - \sqrt[3]{n+1} \right)$.
	\choice 
		{$+\infty$} 
		{$-7$} 
		 {$\dfrac{53}{2}$}
		{\True  $0$}
	\loigiai{ 
		\begin{align*}
			 L& =\lim\left(\sqrt[3]{n+4}-\sqrt[3]{n+1}\right)\\
			&=\lim\dfrac{3}{\sqrt[3]{(n+4)^2}+\sqrt[3]{(n+4) (n+1)}+\sqrt[3]{(n+1)^2}}  \\
			&= \lim\dfrac{3}{\sqrt[3]{n^2 \left(1+\dfrac{4}{n}\right)^2}+\sqrt[3]{n^2 \left(1+\dfrac{4}{n}\right) \left(1+\dfrac{1}{n}\right)}-\sqrt[3]{n^2 \left(1+\dfrac{1}{n}\right)^2}}  \\
			&= \lim\dfrac{3}{\sqrt[3]{n^2}}\left[\sqrt[3]{\left(1+\dfrac{4}{n}\right)^2}+\sqrt[3]{\left(1+\dfrac{4}{n}\right) \left(1+\dfrac{1}{n}\right)}+\sqrt[3]{\left(1+\dfrac{1}{n}\right)^2}\right]\\ 
			& =  0
		\end{align*}
	} 
\end{ex}

% nguồn: Ok_Toan11_TracNghiem_Lan2
\begin{ex}%[1D3N1-3]%[ma:d3743214]
	 Tính giới hạn $L=\lim\left(\sqrt[3]{2n-n^3}+n-1\right)$.
	\choice 
		 {$+\infty$}
		{\True  $-1$}
		 {$\dfrac{53}{2}$}
		{$\dfrac{1}{2}$} 
	 \loigiai{
		\begin{align*}
			L& =\lim\left(\sqrt[3]{2  n-n^3}+n-1\right)\\
			&=-1+\lim\left(\sqrt[3]{2  n-n^3}+n\right)\\
			&=-1+\lim\dfrac{2 n}{\sqrt[3]{\left(2 n-n^3\right)^2}-n  \sqrt[3]{2 n-2 n^3}+n^2} \\
			&  =-1+\lim\dfrac{\dfrac{2}{n}}{\sqrt[3]{\left(\dfrac{2}{n^2}-1\right)^2}-\sqrt[3]{\dfrac{2}{n^2}-1}+1}\\
			&=-1+0\\ 
			 &=-1.
		\end{align*}
	} 
\end{ex}

% nguồn: Ok_Toan11_TracNghiem_Lan2
\begin{ex}%[1D3N1-3]%[ma:c2e9dce1]
	Tính giới hạn  $L=\lim\left(\sqrt[3]{n-n^3}+n+2\right)$.
	\choice 
		{$+\infty$} 
		{\True $2$} 
		{$1$} 
		{$\dfrac{1}{2}$} 
	 \loigiai{
		\begin{align*}
			L&=\lim\left(\sqrt[3]{n-n^3}+n+2\right)\\ 
			&=2+\lim\left(\sqrt[3]{n-n^3}+n\right)\\ 
			 &=2+\lim\dfrac{n}{\sqrt[3]{\left(n-n^3\right)^2}-n \sqrt[3]{n-n^3}+n^2} \\
			 &=2+\lim\dfrac{\dfrac{1}{n}}{\sqrt[3]{\left(\dfrac{1}{n^2}-1\right)^2}-\sqrt[3]{\dfrac{1}{n^2}-1}+1}\\
			 &=2+0\\
			&=2. 
		\end{align*}
	 }
\end{ex}

% nguồn: Ok_Toan11_TracNghiem_Lan2
\begin{ex}%[1D3N1-3]%[ma:5cbbb3c1]
	Tính giới  hạn $L=\lim\left(\sqrt{n^2+3n+5}-n+25\right)$.
	\choice 
		 {$+\infty$}
		{$-7$} 
		{\True  $\dfrac{53}{2}$}
		 {$\dfrac{9}{4}$}
	 \loigiai{
		\begin{align*}
			L&=\lim 25+\lim\left(\sqrt{n^2+3 n+5}-n\right)\\ 
			 &=25+\lim\dfrac{\left(\sqrt{n^2+3 n+5}-n\right)\left(\sqrt{n^2+3 n+5}+n\right)}{\sqrt{n^2+3 n+5}+n}\\
			&=25+\lim\dfrac{\left(n^2+3 n+5\right)-n^2}{\sqrt{n^2+3 n+5}+n}\\ 
			&=25+\lim\dfrac{3  n+5}{\sqrt{n^2+3 n+5}+n} \\
			 &=25+\lim\dfrac{n\left(3+\dfrac{5}{n}\right)}{n\left(\sqrt{1+\dfrac{3}{n}+\dfrac{5}{n^2}}+1\right)}\\
			&=25+\lim\dfrac{3+\dfrac{5}{n}}{\sqrt{1+\dfrac{3}{n}+\dfrac{5}{n^2}}+1}\\ 
			&=25+\dfrac{3+0}{\sqrt{1+0+0}+1}\\ 
			 &=\dfrac{53}{2}.
		\end{align*}
	} 
\end{ex}

% nguồn: Ok_Toan11_TracNghiem_Lan2
\begin{ex}%[1D3N1-3]%[ma:66cf17bc]
	Tính  $L = \lim\limits_{x\to 1} \dfrac{1-x}{\sqrt{2-x}-1}$.
	\choice 
		{$L  = -6$}
		{$L =  -4$}
		{\True $L =  2$}
		{$L =  -2$}
	 \loigiai{
		$L = \lim\limits_{x\to 1}  \dfrac{1-x}{\sqrt{2-x}-1}
		 = \lim\limits_{x\to 1} \dfrac{(1-x)(\sqrt{2-x}+1)}{-x+1}
		= \lim\limits_{x\to  1} (\sqrt{2-x}+1) = 2$.
	 }
\end{ex}

% nguồn: Ok_Toan11_TracNghiem_Lan2
\begin{ex}%[1D3N1-3]%[ma:b60fd52d]
	Cho dãy số $u_n=n\left(\sqrt{n^2+1}-n\right)$. Khi đó  $\lim u_n$ bằng
	\choice 
		{$+\infty$} 
		{$1$} 
		{$0$} 
		{\True $\dfrac{1}{2}$} 
	\loigiai{ 
		 Ta có
		\begin{align*}
			\lim  u_n
			&=\lim  n\left(\sqrt{n^2+1}-n\right)\\
			 &=\lim\dfrac{n\left(\sqrt{n^2+1}-n\right)\left(\sqrt{n^2+1}+n\right)}{\sqrt{n^2+1}+n}\\
			&=\lim\dfrac{n}{\sqrt{n^2+1}+n}\\ 
			 &=\lim\dfrac{1}{\sqrt{1+\dfrac{1}{n^2}}+1}\\
			&=\dfrac{1}{2}. 
		\end{align*}
		Vậy $\lim  u_n=\dfrac{1}{2}$.
	 }
\end{ex}

% nguồn: Ok_Toan11_TracNghiem_Lan2
\begin{ex}%[1D3N1-3]%[ma:1b1296b0]
	Tính $\lim\limits_{x\to  \sqrt{3}} \dfrac{2x^2-6}{x-\sqrt{3}} = a\sqrt{b}$ ($a$ là số nguyên, $b$ là số nguyên tố). Khi đó giá trị của $a+b$ bằng
	\choice 
		 {\True $7$}
		{$10$} 
		 {$5$}
		 {$6$}
	 \loigiai{
		Ta  có
		$\lim\limits_{x\to \sqrt{3}} \dfrac{2x^2-6}{x-\sqrt{3}} 
		= \lim\limits_{x\to \sqrt{3}}  \dfrac{2(x^2-3)}{x-\sqrt{3}}
		= \lim\limits_{x\to \sqrt{3}} 2(x+\sqrt{3}) 
		= 4\sqrt{3}$.\\ 
		Vậy $a+b=7$. 
	 }
\end{ex}

% nguồn: Ok_Toan11_TracNghiem_Lan2
\begin{ex}%[1D3N1-3]%[ma:9be8ff9c]
	 Tính giới hạn $L=\lim\left(\sqrt[3]{8n^3+3n^2+4}-2n+6\right)$.
	 \choice
		 {$+\infty$}
		 {\True $\dfrac{25}{4}$}
		 {$\dfrac{53}{2}$}
		{$\dfrac{1}{2}$} 
	 \loigiai{
		\begin{align*}
			L& =\lim\left(\sqrt[3]{8 n^3+3 n^2+4}-2  n+6\right)\\
			&=6+\lim\left(\sqrt[3]{8 n^3+3 n^2+4}-2 n\right) \\ 
			& =6+\lim\dfrac{3 n^2+4}{\sqrt[3]{\left(8  n^3+3 n^2+4\right)^2}+2 n \sqrt[3]{8 n^3+3 n^2+4}+4 n^2}\\
			 & =6+\lim\dfrac{3+\dfrac{4}{n^2}}{\sqrt[3]{\left(8+\dfrac{3}{n}+\dfrac{4}{n^3}\right)^2}+2 \sqrt[3]{8+\dfrac{3}{n}+\dfrac{4}{n^3}}+4}\\
			 &=6+\dfrac{1}{4}\\
			&=\dfrac{25}{4}. 
		\end{align*}
	 }
\end{ex}

% nguồn: Ok_Toan11_TracNghiem_Lan2
\begin{ex}%[1D3N1-3]%[ma:ec03e7c3]
	 Giá trị $\lim\limits_{x\to-\infty} \dfrac{\sqrt{x^2-x}-\sqrt{4x^2+1}}{2x+3}$ bằng
	 \choice
		{$-\dfrac{1}{2}$} 
		{$+\infty$} 
		 {$-\infty$}
		 {\True $\dfrac{1}{2}$}
	 \loigiai{
		Ta có 
		$\lim\limits_{x\to-\infty} \dfrac{\sqrt{x^2-x}-\sqrt{4x^2+1}}{2x+3} 
		= \lim\limits_{x\to-\infty} \dfrac{-\sqrt{1-\frac{1}{x}}+\sqrt{4+\frac{1}{x^2}}}{2+\frac{3}{x}} 
		=  \dfrac{-\sqrt{1}+\sqrt{4}}{2}
		= \dfrac{1}{2}$. 
	 }
\end{ex}

% nguồn: Ok_Toan11_TracNghiem_Lan2
\begin{ex}%[1D3N1-3]%[ma:a5245af2]
	Tính $I=\lim\left[n\left(\sqrt{n^2+2}-\sqrt{n^2-1}  \right)\right]$.
	\choice 
		 {$I=+\infty$}
		 {\True $I=\dfrac{3}{2}$}
		{$I=1{,}499$} 
		 {$I=0$}
	 \loigiai{
		 Ta có
		\begin{align*}
			I&=\lim\left[n\left(\sqrt{n^2+2}-\sqrt{n^2-1} \right)\right]\\ 
			 &=\lim\dfrac{n\left(\sqrt{n^2+2}-\sqrt{n^2-1} \right)\left(\sqrt{n^2+2}+\sqrt{n^2-1} \right)}{\sqrt{n^2+2}+\sqrt{n^2-1}}\\
			&=\lim\dfrac{3n}{\sqrt{n^2+2}+\sqrt{n^2-1}}\\ 
			&=\lim\dfrac{3}{\sqrt{1+\dfrac{2}{n^2}}+\sqrt{1-\dfrac{1}{n^2}}}\\ 
			&=\dfrac{3}{2}. 
		\end{align*}
	 }
\end{ex}

% nguồn: Ok_Toan11_TracNghiem_Lan2
\begin{ex}%[1D3N1-3]%[ma:37322168]
	 Tính giới hạn $L=\lim\dfrac{\sqrt{2n+1}-\sqrt{n+3}}{\sqrt{4n-5}}$.
	 \choice
		 {$+\infty$}
		{$-7$} 
		{$\dfrac{53}{2}$} 
		 {\True $\dfrac{\sqrt{2}-1}{2}$}
	\loigiai{ 
		\begin{align*}
			L&=\lim\dfrac{\left(\sqrt{2 n+1}-\sqrt{n+3}\right)\left(\sqrt{2  n+1}+\sqrt{n+3}\right)}{\sqrt{4 n-5}\left(\sqrt{2 n+1}+\sqrt{n+3}\right)}\\
			&=\lim\dfrac{(2 n+1)-(n+3)}{\sqrt{4 n-5}\left(\sqrt{2  n+1}+\sqrt{n+3}\right)}\\
			 &=\lim\dfrac{n-2}{\sqrt{4 n-5}\left(\sqrt{2 n+1}+\sqrt{n+3}\right)}\\
			&=\lim\dfrac{n\left(1-\dfrac{2}{n}\right)}{n \sqrt{4-\dfrac{5}{n}}\left(\sqrt{2+\dfrac{1}{n}}+\sqrt{1+\dfrac{3}{n}}\right)} \\ 
			& =\lim\dfrac{1-\dfrac{2}{n}}{\sqrt{4-\dfrac{5}{n}}\left(\sqrt{2+\dfrac{1}{n}}+\sqrt{1+\dfrac{3}{n}}\right)}\\ 
			 &=\dfrac{1-0}{\sqrt{4-0}\left(\sqrt{2+0}+\sqrt{1+0}\right)}\\
			&=\dfrac{\sqrt{2}-1}{2}  .
		\end{align*}
	 }
\end{ex}

% nguồn: Ok_Toan11_TracNghiem_Lan2
\begin{ex}%[1D3N1-3]%[ma:4c074629]
	Tính $I = \lim\limits_{x\to 1} \dfrac{2x-\sqrt{x+3}}{x^2-1}$. 
	\choice 
		{\True $I  = \dfrac{7}{8}$}
		 {$I = \dfrac{3}{2}$}
		 {$I = \dfrac{3}{8}$}
		 {$I = \dfrac{3}{4}$}
	 \loigiai{
		Ta  có
		$I = \lim\limits_{x\to 1}  \dfrac{2x-\sqrt{x+3}}{x^2-1}
		= \lim\limits_{x\to 1} \dfrac{(2x-\sqrt{x+3})(2x+\sqrt{x+3})}{(x-1)(x+1)(2x+\sqrt{x+3})}$.\\ 
		 Rút gọn
		 $I
		= \lim\limits_{x\to 1} \dfrac{4x^2-x-3}{(x-1)(x+1)(2x+\sqrt{x+3})} 
		= \lim\limits_{x\to  1} \dfrac{4x+3}{(x+1)(2x+\sqrt{x+3})}
		=  \dfrac{7}{8}$.
	} 
\end{ex}

% nguồn: Ok_TracNghiemLop11_Lan1
\begin{ex}%[1D3N1-3]%[ma:622cea9d]
	Cho hai dãy $\left(u_n \right)$  và $\left(v_n \right)$ thỏa mãn $\displaystyle\lim\limits u_n = 6$ và $\displaystyle\lim\limits v_n = -2$. Giá trị của $\displaystyle\lim\limits \left(u_n - v_n \right)$ bằng
	\choice 
		{\True  $8$}
		 {$4$}
		{$-6$} 
		{$6$} 
	\loigiai{ 
		Ta có $\displaystyle\lim\limits \left(u_n - v_n\right) = \displaystyle\lim\limits  u_n -\displaystyle\lim\limits v_n = 6 - \left(-2\right)=8$.
	 }
\end{ex}

% nguồn: Ok_TracNghiemLop11_Lan1
\begin{ex}%[1D3N1-3]%[ma:e4676088]
	Giá trị  của giới hạn $\lim \left(\sqrt{{n^2}+2n}-\sqrt{{n^2}-2n} \right)$ là
	 \choice
		 {$1$}
		 {\True $2$}
		{$4$} 
		 {$+\infty$}
	\loigiai{ 
		Vì $\sqrt{{n^2}+2n}-\sqrt{{n^2}-2n}\sim \sqrt{{n^2}}-\sqrt{{n^2}}=0$ nên ta sử dụng phương pháp nhân  lượng liên hợp:
		\[\lim \left(\sqrt{{n^2}+2n}-\sqrt{{n^2}-2n} \right)=\lim \dfrac{4n}{\sqrt{{n^2}+2n}+\sqrt{{n^2}-2n}}=\lim  \dfrac{4}{\sqrt{1+\dfrac{2}{n}}+\sqrt{1-\dfrac{2}{n}}}=2.\]
		Giải nhanh: 
		\[\sqrt{{n^2}+2n}-\sqrt{{n^2}-2n}=\dfrac{4n}{\sqrt{{n^2}+2n}+\sqrt{{n^2}-2n}}\sim \dfrac{4n}{\sqrt{{n^2}}+\sqrt{{n^2}}}=2.\] 
	 }
\end{ex}

% nguồn: Ok_TracNghiemLop11_Lan1
\begin{ex}%[1D3N1-3]%[ma:4f688892]
	Giá trị  của giới hạn $\lim \left(\sqrt{2{n^2}-n+1}-\sqrt{2{n^2}-3n+2} \right)$ là
	 \choice
		 {$0$}
		 {\True $\dfrac{\sqrt{2}}{2}$}
		 {$-\infty$}
		{$+\infty$} 
	 \loigiai{
		Vì $\sqrt{2{n^2}-n+1}-\sqrt{2{n^2}-3n+2}\sim \sqrt{2{n^2}}-\sqrt{2{n^2}}=0$ nên ta sử dụng phương pháp nhân lượng  liên hợp:
		$$\begin{aligned}
		& \lim \left(\sqrt{2{n^2}-n+1}-\sqrt{2{n^2}-3n+2} \right)=\lim \dfrac{2n-1}{\sqrt{2{n^2}-n+1}+\sqrt{2{n^2}-3n+2}} \\ &=\lim  \dfrac{2-\dfrac{1}{n}}{\sqrt{2-\dfrac{1}{n}+\dfrac{1}{{n^2}}}+\sqrt{2-\dfrac{3}{n}+\dfrac{2}{{n^2}}}}=\dfrac{1}{\sqrt{2}}.
	\end{aligned}$$
	Giải nhanh: 
	$$\sqrt{2{n^2}-n+1}-\sqrt{2{n^2}-3n+2}=\dfrac{2n-1}{\sqrt{2{n^2}-n+1}+\sqrt{2{n^2}-3n+2}}\sim  \dfrac{2n}{\sqrt{2{n^2}}+\sqrt{2{n^2}}}=\dfrac{1}{\sqrt{2}}.$$
} 
\end{ex}

% nguồn: Ok_Toan11_TracNghiem_Lan2
\begin{ex}%[1D3H1-3]%[ma:6115becd]
	Biết $\lim\left(n+1-\sqrt{n^2+n}\right)=\dfrac{a}{b}$  với $a,b$ là các số nguyên dương và $\dfrac{a}{b}$ là phân số tối giản. Khi đó $a+2b$ bằng
	 \choice
		{$3$} 
		 {$1$}
		{\True $5$} 
		 {$2$}
	\loigiai{Ta có \allowdisplaybreaks\begin{eqnarray*}
		 \lim\left(n+1-\sqrt{n^2+n}\right)&=&\lim\dfrac{(n+1)^2-\left(n^2+n\right)}{n+1+\sqrt{n^2-n}}\\
		&=&\lim\dfrac{n+1}{n+1+\sqrt{n^2+n}}\\ 
		 &=&\lim\dfrac{1+\dfrac{1}{n}}{1+\dfrac{1}{n}+\sqrt{1+\dfrac{1}{n}}}=\dfrac{1}{2}.
	\end{eqnarray*}
	Suy  ra $a=1$, $b=2$.\\
	 Vậy $a+2b=5$.
 }
\end{ex}

% nguồn: Ok_Toan11_TracNghiem_Lan2
\begin{ex}%[1D3H1-3]%[ma:258319cf]
	Cho dãy số $(u_n)$ với $u_n=\sqrt{n^2+a n+5}-\sqrt{n^2+1}$, trong đó $a$ là tham số thực. Tìm $a$  để $\lim u_n=-1$.
	\choice 
		{$3$} 
		 {$2$}
		 {\True $-2$}
		{$-3$} 
	 \loigiai{Ta có
		\begin{eqnarray*}
			\lim u_n&=&\lim\left(\sqrt{n^2+a n+5}-\sqrt{n^2+1}\right)\\ 
			&=&\lim\dfrac{a n+4}{\sqrt{n^2+a  n+5}+\sqrt{n^2+1}}\\
			 &=&\lim\dfrac{a+\dfrac{4}{n}}{\sqrt{1+\dfrac{a}{n}+\dfrac{5}{n^2}}+\sqrt{1+\dfrac{1}{n^2}}}\\
			 &=&\dfrac{a}{2}.
		\end{eqnarray*}
		Mà $\lim  u_n=-1$ nên $\dfrac{a}{2}=-1$ hay $a=-2$.
	 }
\end{ex}

% nguồn: Ok_Toan11_TracNghiem_Lan2
\begin{ex}%[1D3H1-3]%[ma:415cae5e]
	Có bao nhiêu giá trị nguyên của $a$ để $\lim\left(\sqrt{n^2-4n+7}+a-n\right)=0$? 
	 \choice
		 {$3$}
		 {$1$}
		{\True $2$} 
		{$0$} 
	 \loigiai{
		Ta có 
		\begin{align*}
			\lim\left(\sqrt{n^2-4n+7}+a-n\right) 
			 &=\lim\dfrac{\left(\sqrt{n^2-4n+7}+a-n\right)\left[\sqrt{n^2-4n+7}-\left(a-n\right)\right]}{\sqrt{n^2-4n+7}-\left(a-n\right)}\\
			&=\lim\dfrac{-4n+7+2an-a^2}{\sqrt{n^2-4n+7}-\left(a-n\right)}\\ 
			 &=\lim\dfrac{2a-4+\dfrac{7-a^2}{n}}{\sqrt{1-\dfrac{4}{n}+\dfrac{7}{n^2}}-\dfrac{a}{n}+1}\\
			&=a-2. 
		\end{align*}
		Để $\lim\left(\sqrt{n^2-4n+7}+a-n\right)=0$ thì  $a-2=0\Leftrightarrow a=2$.
	} 
\end{ex}

% nguồn: Ok_Toan11_TracNghiem_Lan2
\begin{ex}%[1D3H1-3]%[ma:a5ebf5f5]
	Cho $\lim\left(\sqrt{n^2+an}-n\right)=\dfrac{\sqrt{2}}{2}$. Số thực $a$ thuộc khoảng nào  sau đây?
	 \choice
		 {$(2;3)$}
		{\True  $(1;2)$}
		{$(-1;0)$} 
		 {$(0;1)$}
	\loigiai{Ta có $\lim\left(\sqrt{n^2+an}-n\right)=\lim\dfrac{an}{\sqrt{n^2+an}+n}=\lim\dfrac{a}{\sqrt{1+\dfrac{a}{n}}+1}=\dfrac{a}{2}$.\\ 
		Theo đề bài thì  $\dfrac{a}{2}=\dfrac{\sqrt{2}}{2}\Leftrightarrow a=\sqrt{2}\in(1;2)$.
	 }
\end{ex}

% nguồn: Ok_Toan11_TracNghiem_Lan2
\begin{ex}%[1D3H1-3]%[ma:d2b6f247]
	Giới hạn $\lim\left(\sqrt{n^2+5n+1}-n\right)$  bằng
	 \choice
		 {$\dfrac{3}{2}$}
		{$0$} 
		 {$\dfrac{1}{4}$}
		{\True  $\dfrac{5}{2}$}
	 \loigiai{
		Ta  có
		$$ 
		\lim\left(\sqrt{n^2+5n+1}-n\right)=\lim\dfrac{5n+1}{\sqrt{n^2+5n+1}+n}=\lim\dfrac{5+\dfrac{1}{n}}{\sqrt{1+\dfrac{5}{n}+\dfrac{1}{n^2}}+1}=\dfrac{5}{2}. 
	$$} 
\end{ex}

% nguồn: Ok_Toan11_TracNghiem_Lan2
\begin{ex}%[1D3H1-3]%[ma:5402c561]
	Tính  $\lim\left(n-\sqrt{n^2-4n+1}\right)$.
	 \choice
		{$-\infty$} 
		 {$-2$}
		 {$+\infty$}
		{\True  $2$}
	\loigiai{ 
		Ta có \allowdisplaybreaks\begin{eqnarray*}
		 \lim\left(n-\sqrt{n^2-4n+1}\right)&=&\lim\dfrac{4n-1}{n+\sqrt{n^2-4n+1}}\\
		 &=&\lim\dfrac{4-\dfrac{1}{n}}{1+\sqrt{1-\dfrac{4}{n}+\dfrac{1}{n^2}}}=\dfrac{4-0}{1+\sqrt{1-0+0}}=2.
	\end{eqnarray*}
 }
\end{ex}

% nguồn: Ok_Toan11_TracNghiem_Lan2
\begin{ex}%[1D3H1-3]%[ma:c78333b7]
	Cho dãy số $\left(u_n \right)$ với $u_n=\sqrt{n^2+an-3}-\sqrt{n^2+n}$, trong đó $a$ là tham  số thực. Tìm $a$ để $\lim u_n=3$.
	\choice 
		{\True $7$} 
		 {$4$}
		{$5$} 
		 {$6$}
	 \loigiai{
		Ta có 
		\begin{align*}
			\lim  u_n=3
			&\Leftrightarrow \lim\left(\sqrt{n^2+an-3}-\sqrt{n^2+n}  \right)=3\\
			&\Leftrightarrow \lim\dfrac{\left(\sqrt{n^2+an-3}-\sqrt{n^2+n} \right)\left(\sqrt{n^2+an-3}+\sqrt{n^2+n}  \right)}{\sqrt{n^2+an-3}+\sqrt{n^2+n}}=3\\
			&\Leftrightarrow  \lim\dfrac{\left(a-1\right)n-3}{\sqrt{n^2+an-3}+\sqrt{n^2+n}}=3\\
			 &\Leftrightarrow \lim\dfrac{\left(a-1\right)-\dfrac{3}{n}}{\sqrt{1+\dfrac{a}{n}-\dfrac{3}{n^2}}+\sqrt{1+\dfrac{1}{n}}}=3\\
			&\Leftrightarrow  \dfrac{a-1}{2}=3\\
			&\Leftrightarrow  a=7.
		\end{align*}
		Vậy giá  trị của $a$ cần tìm là $a=7$.
	} 
\end{ex}

% nguồn: Ok_Toan11_TracNghiem_Lan2
\begin{ex}%[1D3H1-3]%[ma:8a4141a8]
	Giá trị của giới  hạn $\lim\left(\sqrt{3n+5}-\sqrt{3n+1}\right)$ bằng
	 \choice
		{\True $0$} 
		{$1$} 
		{$3$} 
		 {$5$}
	\loigiai{ 
		 Ta có
		\allowdisplaybreaks\begin{eqnarray*}
		\lim\left(\sqrt{3n+5}-\sqrt{3n+1}\right)&=&\lim\dfrac{\left(\sqrt{3n+5}-\sqrt{3n+1}\right)\left(\sqrt{3n+5}+\sqrt{3n+1}\right)}{\sqrt{3n+5}+\sqrt{3n+1}}\\ 
		&=&\lim\dfrac{4}{\sqrt{3n+5}+\sqrt{3n+1}}\\ 
		&=&\lim\dfrac{\dfrac{4}{\sqrt{n}}}{\sqrt{3+\dfrac{5}{n}}+\sqrt{3+\dfrac{1}{n}}}=0. 
	\end{eqnarray*}
 }
\end{ex}

% nguồn: Ok_Toan11_TracNghiem_Lan2
\begin{ex}%[1D3H1-3]%[ma:c3f0e28e]
	Giá trị đúng của  $\lim\sqrt{n}\left(\sqrt{4n+1}-\sqrt{4n-1}\right)$ là
	\choice 
		 {$0$}
		 {\True $\dfrac{1}{2}$}
		 {$+\infty$}
		 {$\dfrac{-1}{2}$}
	\loigiai{ 
		Ta có \allowdisplaybreaks\begin{eqnarray*}
		\lim\sqrt{n}\left(\sqrt{4n+1}-\sqrt{4n-1}\right)&=&\lim\dfrac{2\sqrt{n}}{\sqrt{4n+1}+\sqrt{4n-1}}\\ 
		 &=&\lim\dfrac{2}{\sqrt{4+\dfrac{1}{n}}+\sqrt{4-\dfrac{1}{n}}}=\dfrac{1}{2}.
	\end{eqnarray*}
 }
\end{ex}

% nguồn: Ok_TracNghiemLop11_Lan1
\begin{ex}%[1D3H1-3]%[ma:2573ec9e]
	Giá trị của giới hạn $\lim \left(\sqrt[3]{{n^3}-2{n^2}}-n  \right)$ bằng
	 \choice
		 {$\dfrac{1}{3}$}
		 {\True $-\dfrac{2}{3}$}
		 {$0$}
		 {$1$}
	\loigiai{ 
		Sử  dụng phương pháp nhân lượng liên hợp:
		$$ 
		\begin{aligned}
			\lim\left(\sqrt[3]{{n^3}-2{n^2}}-n\right)&=\lim\dfrac{-2n^2}{\sqrt[3]{\left({n^3}-2{n^2}\right)^2}+n \sqrt[3]{{n^3}-2{n^2}}+n^2 }\\ 
			&=\lim\dfrac{-2}{\sqrt[3]{\left(1-\dfrac{2}{n} \right)^2 }+\sqrt[3]{1-\dfrac{2}{n}}+1}=-\dfrac{2}{3}. 
		\end{aligned}
		$$ 
		 Giải nhanh:
		\[\sqrt[3]{{n^3}-2{n^2}}-n=\dfrac{-2{n^2}}{\sqrt[3]{{{\left({n^3}-2{n^2} \right)}^2}}+n.\sqrt[3]{{n^3}-2{n^2}}+{n^2}}\sim \dfrac{-2{n^2}}{\sqrt[3]{{n^6}}+n.\sqrt[3]{{n^3}}+{n^2}}=-\dfrac{2}{3}.\] 
	 }
\end{ex}

% nguồn: Ok_TracNghiemLop11_Lan1
\begin{ex}%[1D3H1-3]%[ma:eb9ed498]
	Có bao nhiêu giá trị nguyên của $a$ thỏa $\lim \left(\sqrt{{n^2}-8n}-n+{a^2}  \right)=0.$
	 \choice
		{$0$} 
		{\True  $2$}
		 {$1$}
		 {Vô số}
	 \loigiai{
		Vì $\sqrt{{n^2}-8n}-n+{a^2}\sim \sqrt{{n^2}}-n=0$ nên ta sử  dụng phương pháp nhân lượng liên hợp:
		\[\lim \left(\sqrt{{n^2}-8n}-n+{a^2} \right)=\lim  \dfrac{\left(2{a^2}-8\right)n}{\sqrt{{n^2}+n}+n}=\lim \dfrac{2{a^2}-8}{\sqrt{1+\dfrac{1}{n}}+1}={a^2}-4=0.\]
		Do đó $ a=\pm  2$.\\
		Vậy có $2$ giá trị  nguyên của $a$ thỏa mãn yêu cầu bài toán.
	} 
\end{ex}

% nguồn: Ok_TracNghiemLop11_Lan1
\begin{ex}%[1D3H1-3]%[ma:6c59794a]
	Giá trị của giới hạn $\lim \left(\sqrt{{n^2}+2n-1}-\sqrt{2{n^2}+n} \right)$  là
	\choice 
		{$-1$} 
		 {$1-\sqrt{2}$}
		{\True $-\infty$} 
		 {$+\infty$}
	\loigiai{ 
		 Ta có $\lim \left(\sqrt{{n^2}+2n-1}-\sqrt{2{n^2}+n} \right)=\lim n.\left(\sqrt{1+\dfrac{2}{n}-\dfrac{1}{{n^2}}}-\sqrt{2+\dfrac{1}{n}} \right)=-\infty$ vì
		$\lim n=+\infty, \lim \left(\sqrt{1+\dfrac{2}{n}-\dfrac{1}{{n^2}}}-\sqrt{2+\dfrac{1}{n}} \right)=1-\sqrt{2} <  0$.\\
		 Giải nhanh: $\sqrt{{n^2}+2n-1}-\sqrt{2{n^2}+n}\sim \sqrt{{n^2}}-\sqrt{2{n^2}}=\left(1-\sqrt{2} \right)n\xrightarrow[{}]{}-\infty$.
	} 
\end{ex}

% nguồn: Ok_TracNghiemLop11_Lan1
\begin{ex}%[1D3H1-3]%[ma:e154fa80]
	Có bao nhiêu  giá trị của $a$ để $\lim \left(\sqrt{{n^2}+{a^2}n}-\sqrt{{n^2}+\left(a+2\right)n+1} \right)=0$.
	\choice 
		{$0$} 
		{\True  $2$}
		 {$1$}
		 {$3$}
	 \loigiai{
		Vì $\sqrt{{n^2}+{a^2}n}-\sqrt{{n^2}+\left(a+2\right)n+1}\sim \sqrt{{n^2}}-\sqrt{{n^2}}=0$ nên ta  sử dụng phương pháp nhân lượng liên hợp:
		$$\begin{aligned}
		&\lim  \left(\sqrt{{n^2}+{a^2}n}-\sqrt{{n^2}+\left(a+2\right)n+1} \right)=\lim \dfrac{\left({a^2}-a-2\right)n-1}{\sqrt{{n^2}+n}+\sqrt{{n^2}+1}}\\
		&=\lim \dfrac{{a^2}-a-2-\dfrac{1}{n}}{\sqrt{1+\dfrac{1}{n}}+\sqrt{1+\dfrac{1}{{n^2}}}}=\dfrac{{a^2}-a-2}{2}=0\Leftrightarrow  \hoac{& a=-1\\
		& a=2\\ 
		}. 
	\end{aligned}$$
} 
\end{ex}

% nguồn: Ok_TracNghiemLop11_Lan1
\begin{ex}%[1D3H1-3]%[ma:3125dce1]
	Giá  trị của giới hạn $\lim \left(\sqrt{{n^2}-2n+3}-n \right)$ là
	\choice 
		{\True  $-1$}
		 {$0$}
		 {$1$}
		{$+\infty$} 
	 \loigiai{
		Vì $\sqrt{{n^2}-2n+3}-n\sim \sqrt{{n^2}}-n=0$ nên ta  sử dụng phương pháp nhân lượng liên hợp:
		$$\lim \left(\sqrt{{n^2}-2n+3}-n \right)=\lim \dfrac{-2n+3}{\sqrt{{n^2}-2n+3}+n}=\lim \dfrac{-2+\dfrac{3}{n}}{\sqrt{1-\dfrac{2}{n}+\dfrac{3}{{n^2}}}+1}=-1.$$ 
		Giải nhanh: $\sqrt{{n^2}-2n+3}-n=\dfrac{-2n+3}{\sqrt{{n^2}-2n+3}+n}\sim  \dfrac{-2n}{\sqrt{{n^2}}+n}=-1$.
	 }
\end{ex}

% nguồn: Ok_TracNghiemLop11_Lan1
\begin{ex}%[1D3H1-3]%[ma:b4c3369b]
	Giá trị của giới hạn $\lim \left(\sqrt[3]{{n^3}+1}-\sqrt[3]{{n^3}+2}  \right)$
	 bằng
	 \choice
		{$3$} 
		 {$2$}
		 {\True $0$}
		 {$1$}
	\loigiai{ 
		Vì $\sqrt[3]{{n^3}+1}-\sqrt[3]{{n^3}+2}\sim \sqrt[3]{{n^3}}-\sqrt[3]{{n^3}}=0$ nên ta  sử dụng phương pháp nhân lượng liên hợp:
		\[\lim \left(\sqrt[3]{{n^3}+1}-\sqrt[3]{{n^3}+2}  \right)=\lim \dfrac{-1}{\sqrt[3]{{{\left({n^3}+1\right)}^2}}+\sqrt[3]{{n^3}+1}\cdot\sqrt[3]{{n^3}+2}+\sqrt[3]{\left({n^3}+2\right)}}=0.\]
	 }
\end{ex}

% nguồn: Ok_TracNghiemLop11_Lan1
\begin{ex}%[1D3H1-3]%[ma:b895af4b]
	Giá trị của giới hạn  $\lim \dfrac{1}{\sqrt[3]{{n^3}+1}-n}$
	là: 
	 \choice
		{$2$} 
		{\True  $0$}
		 {$-\infty$}
		{$+\infty$} 
	\loigiai{ 
		Vì $\sqrt[3]{{n^3}+1}-n\sim \sqrt[3]{{n^3}}-n=0$ nên ta  sử dụng phương pháp nhân lượng liên hợp:
		\[\lim \left(\sqrt[3]{{n^3}+1}-n \right)=\lim  \dfrac{1}{\sqrt[3]{{{\left({n^3}+1\right)}^2}}+n\sqrt[3]{{n^3}+1}+{n^2}}=0.\]
	} 
\end{ex}

% nguồn: Ok_TracNghiemLop11_Lan1
\begin{ex}%[1D3H1-3]%[ma:32fa4096]
	 Cho dãy số $\left({u_n} \right)$ với ${u_n}=\sqrt{{n^2}+an+5}-\sqrt{{n^2}+1}$, trong đó $a$ là tham số thực. Tìm $a$ để $\lim {u_n}=-1$.
	\choice 
		{$3$} 
		{$2$} 
		{\True  $-2$}
		 {$-3$}
	\loigiai{ 
		Vì $\sqrt{{n^2}+an+5}-\sqrt{{n^2}+1}\sim \sqrt{{n^2}}-\sqrt{{n^2}}=0$ nên ta sử dụng  phương pháp nhân lượng liên hợp:
		\[\begin{aligned}
		-1&=\lim {u_n}=\lim \left(\sqrt{{n^2}+an+5}-\sqrt{{n^2}+1} \right)=\lim  \dfrac{an+4}{\sqrt{{n^2}+an+5}+\sqrt{{n^2}+1}} \\
		&=\lim  \dfrac{a+\dfrac{4}{n}}{\sqrt{1+\dfrac{a}{n}+\dfrac{5}{{n^2}}}+\sqrt{1+\dfrac{1}{{n^2}}}}=\dfrac{a}{2}\Leftrightarrow a=-2.
	\end{aligned}\]
	 Giải nhanh:
	\[-1\sim \sqrt{{n^2}+an+5}-\sqrt{{n^2}+1}=\dfrac{an+4}{\sqrt{{n^2}+an+5}+\sqrt{{n^2}+1}}\sim \dfrac{an}{\sqrt{{n^2}}+\sqrt{{n^2}}}=\dfrac{a}{2}\Leftrightarrow a=-2.\] 
} 
\end{ex}

% nguồn: Ok_TracNghiemLop11_Lan1
\begin{ex}%[1D3H1-3]%[ma:3c6c0977]
	Giá trị của giới hạn $\lim \dfrac{\sqrt{9{n^2}-n}-\sqrt{n+2}}{3n-2}$  là
	\choice 
		{\True $1$} 
		 {$0$}
		 {$3$}
		 {$+\infty$}
	\loigiai{ 
		Vì $\sqrt{9{n^2}-n}-\sqrt{n+2}\sim \sqrt{9{n^2}}=3n \ne 0$ nên chia cả tử và mẫu với $n$ ta  có
		 $$\lim \dfrac{\sqrt{9{n^2}-n}-\sqrt{n+2}}{3n-2}=\lim \dfrac{\sqrt{9-\dfrac{1}{n}}-\sqrt{\dfrac{1}{n}+\dfrac{2}{{n^2}}}}{3-\dfrac{2}{n}}=\dfrac{\sqrt{9}}{3}=1.$$ Giải nhanh: $\dfrac{\sqrt{9{n^2}-n}-\sqrt{n+2}}{3n-2}\sim \dfrac{\sqrt{9{n^2}}}{3n}=1$.
	 }
\end{ex}

% nguồn: Ok_Toan11_TracNghiem_Lan2
\begin{ex}%[1D3V1-3]%[ma:c1f465d3]
	Tính giới hạn $L=\lim\left(\sqrt{n^{2}+n+1}-\sqrt[3]{n^{3}+n^{2}}\right)$. 
	 \choice
		 {$+\infty$}
		 {$\dfrac{5}{4}$}
		{$\dfrac{53}{2}$} 
		{\True $\dfrac{1}{6}$} 
	\loigiai{ 
		Ta  có
		\begin{eqnarray*}
			L&=&\lim\left(\sqrt{n^{2}+n+1}-\sqrt[3]{n^{3}+n^{2}}\right)\\ 
			 &=&\lim\left[\left(\sqrt{n^{2}+n+1}-n\right)+\left(n-\sqrt[3]{n^{3}+n^{2}}\right)\right] \\
			&=&\lim\left[\dfrac{n^{2}+n+1-n^{2}}{\sqrt{n^{2}+n+1}+n}+\dfrac{n^{3}-\left(n^{3}+n^{2}\right)}{n^{2}+n \sqrt[3]{n^{3}+n^{2}}+\left(\sqrt[3]{n^{3}+n^{2}}\right)^{2}}\right]  \\
			&=&\lim\left[\dfrac{n+1}{\sqrt{n^{2}+n+1}+n}-\dfrac{n^{2}}{n^{2}+n \sqrt[3]{n^{3}+n^{2}}+\left(\sqrt[3]{n^{3}+n^{2}}\right)^{2}}\right] \\ 
			&=&\lim\left\{\dfrac{n\left(1+\dfrac{1}{n}\right)}{n\left(\sqrt{1+\dfrac{1}{n}+\dfrac{1}{n^{2}}}+1\right)}-\dfrac{n^{2}}{n^{2}\left[1+\sqrt[3]{1+\dfrac{1}{n}}+\left(\sqrt[3]{1+\dfrac{1}{n}}\right)^{2}\right]}\right\}  \\
			&=&\lim\left[\dfrac{1+\dfrac{1}{n}}{\sqrt{1+\dfrac{1}{n}+\dfrac{1}{n^{2}}}+1}-\dfrac{1}{1+\sqrt[3]{1+\dfrac{1}{n}}+\left(\sqrt[3]{1+\dfrac{1}{n}}\right)^{2}}\right]\\ 
			&=&\dfrac{1}{2}-\dfrac{1}{3}\\ 
			&=&\dfrac{1}{6}. 
		\end{eqnarray*}
	} 
\end{ex}

% nguồn: Ok_Toan11_TracNghiem_Lan2
\begin{ex}%[1D3V1-3]%[ma:c33a0b3d]
	Tính giới hạn $L=\lim\left(\sqrt[3]{n^{3}-2n^{2}}-n-1\right)$. 
	\choice 
		 {$+\infty$}
		{$\dfrac{5}{4}$} 
		 {$\dfrac{53}{2}$}
		{\True $-\dfrac{5}{3}$} 
	\loigiai{ 
		 Ta có
		\begin{eqnarray*}
			L&=&\lim\left(\sqrt[3]{n^{3}-2n^{2}}-n-1\right)\\ 
			&=&-1+\lim\left(\sqrt[3]{n^{3}-2n^{2}}-n\right)\\ 
			&=&-1+\lim\dfrac{-2n^{2}}{\sqrt[3]{\left(n^{3}-2n^{2}\right)^{2}}+n\cdot \sqrt[3]{2n^{3}-2n^{2}}+n^{2}}\\ 
			 &=&-1+\lim\dfrac{-2}{\sqrt[3]{\left(1-\dfrac{2}{n}\right)^{2}}+\sqrt[3]{1-\dfrac{2}{n}}+1}\\
			 &=&-1-\dfrac{2}{3}=-\dfrac{5}{3}.
		\end{eqnarray*}
	 }
\end{ex}

% nguồn: Ok_Toan11_TracNghiem_Lan2
\begin{ex}%[1D3V1-3]%[ma:890c8cc8]
	Có  bao nhiêu giá trị nguyên của $a$ thỏa $\lim\left(\sqrt{n^2-8n}-n+a^2\right)=0$?
	\choice 
		{$0$} 
		 {\True $2$}
		 {$1$}
		{Vô số} 
	\loigiai{ 
		Ta có \allowdisplaybreaks\begin{eqnarray*}
		 \lim\left(\sqrt{n^2-8n}-n+a^2\right)&=&\lim\dfrac{n^2-8n-\left(n-a^2\right)^2}{\sqrt{n^2-8n}+n-a^2}\\
		&=&\lim\dfrac{\left(2a^2-8\right)n-a^4}{\sqrt{n^2-8n}+n-a^2}\\ 
		&=&\lim\dfrac{2a^2-8-\dfrac{a^4}{n}}{\sqrt{1-\dfrac{8}{n}}+1-\dfrac{a^2}{n}}=a^2-4. 
	\end{eqnarray*}
	Theo đề bài thì  $a^2-4=0\Leftrightarrow a=\pm 2$.\\
	Vậy có $2$ giá trị nguyên của $a$  thỏa mãn bài toán.
 }
\end{ex}

% nguồn: Ok_Toan11_TracNghiem_Lan2
\begin{ex}%[1D3V1-3]%[ma:5a462715]
	Tính  giới hạn $T=\lim\left(\sqrt{16^{n+1}+4^{n}}-\sqrt{16^{n+1}+3^{n}}\right)$.
	\choice 
		{$T=0$} 
		{$T=\dfrac{1}{4}$} 
		 {\True $T=\dfrac{1}{8}$}
		 {$T=\dfrac{1}{16}$}
	 \loigiai{
		Ta có \begin{eqnarray*}
		 T&=&\lim\left(\sqrt{16^{n+1}+4^{n}}-\sqrt{16^{n+1}+3^{n}}\right)\\
		 &=&\lim\dfrac{4^{n}-3^{n}}{\sqrt{16^{n+1}+4^{n}}+\sqrt{16^{n+1}+3^{n}}}\\
		&=&\lim\dfrac{4^{n}-3^{n}}{\sqrt{16\cdot 16^{n}+4^{n}}+\sqrt{16\cdot 16^{n}+3^{n}}}\\ 
		 &=&\lim\dfrac{1-\left(\dfrac{3}{4}\right)^{n}}{\sqrt{16+\left(\dfrac{1}{4}\right)^{n}}+\sqrt{16+\left(\dfrac{3}{4}\right)^{n}}}\\
		&=&\dfrac{1}{4+4}\\ 
		 &=&\dfrac{1}{8}.
	\end{eqnarray*}
 }
\end{ex}

% nguồn: Ok_Toan11_TracNghiem_Lan2
\begin{ex}%[1D3V1-3]%[ma:43985f43]
	Biết giới hạn $\lim\left[n\left(\sqrt{n^2+3}-\sqrt{n^2+2}\right)\right]=\dfrac{a}{b}$ với $a\in\mathbb{N}$, $b\in\mathbb{N}^{*}$ và $\dfrac{a}{b}$ là phân số tối giản. Khi đó, giá trị $a^2-b$ bằng 
	 \choice
		{$3$} 
		{\True $-1$} 
		{$4$} 
		{$5$} 
	\loigiai{Ta có \allowdisplaybreaks\begin{eqnarray*}
		 \lim\left[n\left(\sqrt{n^2+3}-\sqrt{n^2+2}\right)\right]&=&\lim\dfrac{n}{\sqrt{n^2+3}+\sqrt{n^2+2}}\\
		 &&\lim\dfrac{1}{\sqrt{1+\dfrac{3}{n^2}}+\sqrt{1+\dfrac{2}{n^2}}}=\dfrac{1}{2}.
	\end{eqnarray*}
	Suy  ra $a=1$, $b=2$.\\
	Vậy $a^2-b=1^2-2=-1$. 
} 
\end{ex}

% nguồn: Ok_Toan11_TracNghiem_Lan2
\begin{ex}%[1D3V1-3]%[ma:e1e9178e]
	Cho các số thực $a,b,c$ thỏa  mãn $c^2+a=8$ và $\lim\left(\sqrt{an^2+bn}-cn\right)=2$.
	Tính $P=a+b+c$ 
	\choice 
		 {$-24$}
		 {$6$}
		 {$12$}
		{\True $14$} 
	\loigiai{Từ giả thiết $\lim\left(\sqrt{an^2+bn}-cn\right)=2$ suy ra $a>0$, $c>0$.\\ 
		Ta có \allowdisplaybreaks\begin{eqnarray*}
		\lim\left(\sqrt{an^2+bn}-cn\right)=2&\Leftrightarrow&\lim\dfrac{\left(a-c^2\right)n^2+bn}{\sqrt{an^2+bn}+cn}=2\\ 
		&\Leftrightarrow&\heva{&a-c^2=0&&  (1)\\&\dfrac{b}{\sqrt{a}+c}=2.&&(2)}
	\end{eqnarray*}
	Mà $c^2+a=8$ nên từ $(1)$ suy ra $a=c^2=4\Rightarrow  c=2$.\\
	 Thay $c=2$ vào $(2)$ ta được $b=8$.\\
	 Vậy $P=a+b+c=14$.
 }
\end{ex}

% nguồn: Ok_Toan11_TracNghiem_Lan2
\begin{ex}%[1D3V1-3]%[ma:99687ffc]
	Tính $\lim n\left(\sqrt{4n^2+3}-\sqrt[3]{8n^3+n} \right)$. 
	 \choice
		{$+\infty$} 
		 {$1$}
		 {$-\infty$}
		{\True $\dfrac{2}{3}$} 
	 \loigiai{
		 Ta có
		\begin{align*}
			\lim n\left(\sqrt{4  n^2+3}-\sqrt[3]{8 n^3+n}\right)
			&=\lim n\left[\left(\sqrt{4  n^2+3}-2 n\right)+\left(2 n-\sqrt[3]{8 n^3+n}\right)\right]\\
			&=\lim\left[n\left(\sqrt{4  n^2+3}-2 n\right)+n\left(2 n-\sqrt[3]{8 n^3+n}\right)\right]
		\end{align*}
		 Ta có
		\begin{align*}
			\lim n\left(\sqrt{4  n^2+3}-2 n\right)
			 &=\lim\dfrac{n\left(\sqrt{4 n^2+3}-2 n\right)\left(\sqrt{4 n^2+3}+2 n\right)}{\sqrt{4 n^2+3}+2 n}\\
			 &=\lim\dfrac{3 n}{\sqrt{4 n^2+3}+2 n}\\
			 &=\lim\dfrac{3}{\sqrt{4+\dfrac{3}{n^2}}+2}\\
			 &=\dfrac{3}{4}.
		\end{align*}
		 Ta có
		\begin{align*}
			 \lim n\left(2 n-\sqrt[3]{8 n^3+n}\right)
			&=\lim\dfrac{-n^2}{4 n^2+2 n \sqrt[3]{8 n^3+n}+\sqrt[3]{\left(8 n^3+n\right)^2}} \\ 
			&=\lim\dfrac{-1}{4+2  \sqrt[3]{8+\dfrac{1}{n^2}}+\sqrt[3]{\left(8+\dfrac{1}{n^2}\right)^2}}\\
			&=-\dfrac{1}{12}  .
		\end{align*}
		Vậy $\lim n\left(\sqrt{4n^2+3}-\sqrt[3]{8n^3+n} \right)=\dfrac{3}{4}-\dfrac{1}{12}$ $=\dfrac{2}{3}$. 
	 }
\end{ex}

% nguồn: Ok_Toan11_TracNghiem_Lan2
\begin{ex}%[1D3V1-3]%[ma:c18fc8a1]
	Tính giới hạn $L=\lim\left(\sqrt{n^{4}+n^{2}}-\sqrt[3]{n^{6}+1}\right)$. 
	\choice 
		{$+\infty$} 
		 {$\dfrac{5}{4}$}
		 {\True $\dfrac{1}{2}$}
		 {$-\dfrac{5}{3}$}
	\loigiai{ 
		 Ta có
		\begin{eqnarray*}
			L&=&\lim\left(\sqrt{n^{4}+n^{2}}-\sqrt[3]{n^{6}+1}\right)\\ 
			 &=&\lim\left[\left(\sqrt{n^{4}+n^{2}}-n^{2}\right)-\left(\sqrt[3]{n^{6}+1}-n^{2}\right)\right] \\
			&=&\lim\left(\sqrt{n^{4}+n^{2}}-n^{2}\right)-\lim\left(\sqrt[3]{n^{6}+1}-n^{2}\right) \\ 
			 &=&\lim\dfrac{\left(n^{4}+n^{2}-n^{4}\right)}{\sqrt{n^{4}+n^{2}}+n^{2}}-\lim\dfrac{\left(n^{6}+1\right)-n^{6}}{\sqrt[3]{\left(n^{6}+1\right)^{2}}+n^{2} \sqrt[3]{n^{6}+1}+n^{4}} \\
			&=&\lim\dfrac{n^{2}}{\sqrt{n^{4}+n^{2}}+n^{2}}-\lim\dfrac{1}{\sqrt[3]{\left(n^{6}+1\right)^{2}}+n^{2} \sqrt[3]{n^{6}+1}+n^{4}}\\ 
			 &=&\lim\dfrac{1}{\sqrt{1+\dfrac{1}{n^{2}}}+1}+0\\
			 &=&\dfrac{1}{2}.
		\end{eqnarray*}
	 }
\end{ex}

% nguồn: Ok_Toan11_TracNghiem_Lan2
\begin{ex}%[1D3V1-3]%[ma:55787309]
	 Giới hạn $\lim\left(\sqrt{an^2+bn+1}-n\right)=\dfrac{3}{2}$ ($a;b\in\mathbb{Z}$). Khi đó $a^3+b^2$ bằng
	\choice 
		 {$12$}
		{$9$} 
		 {$11$}
		 {\True $10$}
	\loigiai{ 
		Ta có \allowdisplaybreaks\begin{eqnarray*}
		\lim\left(\sqrt{an^2+bn+1}-n\right)&=&\lim\dfrac{\left(\sqrt{an^2+bn+1}-n\right)\left(\sqrt{an^2+bn+1}+n\right)}{\sqrt{an^2+bn+1}+n}\\ 
		 &=&\lim\dfrac{an^2+bn+1-n^2}{\sqrt{an^2+bn+1}+n}=\lim\dfrac{n^2(a-1)+bn+1}{\sqrt{an^2+bn+1}+n}.
	\end{eqnarray*}
	Vì 
	$\lim\left(\sqrt{an^2+bn+1}-n\right)=\dfrac{3}{2}$ nên  $$\lim\dfrac{n^2(a-1)+bn+1}{\sqrt{an^2+bn+1}+n}=\dfrac{3}{2}\Leftrightarrow\heva{&a-1=0\\
	 &\dfrac{b}{\sqrt{a}+1}=\dfrac{3}{2}}\Leftrightarrow\heva{&a=1\\
	&b=3.}$$ 
Khi  đó $a^2+b^2=1^2+3^2=10$.}
\end{ex}

% nguồn: Ok_Toan11_TracNghiem_Lan2
\begin{ex}%[1D3V1-3]%[ma:57550e04]
	Giá  trị của giới hạn $\lim\left(\sqrt{n^2+2^{2023}n}-\sqrt{n^2-2^{2023}n}\right)$ là
	 \choice
		{$2^{2024}$} 
		{\True $2^{2023}$} 
		 {$0$}
		{$+\infty$} 
	\loigiai{Ta  có
		\allowdisplaybreaks\begin{eqnarray*}
		 \lim\left(\sqrt{n^2+2^{2023}n}-\sqrt{n^2-2^{2023}n}\right)&=&\lim\dfrac{2^{2024}n}{\sqrt{n^2+2^{2023}n}+\sqrt{n^2-2^{2023}n}}\\
		&=&\lim\dfrac{2^{2024}}{\sqrt{1+\dfrac{2^{2023}}{n}}+\sqrt{1-\dfrac{2^{2023}}{n}}}\\ 
		 &=&2^{2023}.
	\end{eqnarray*}
} 
\end{ex}

% nguồn: Ok_Toan11_TracNghiem_Lan2
\begin{ex}%[1D3V1-3]%[ma:d108147f]
	Cho $a$, $b$  là các số thực thỏa mãn $\lim\left(\sqrt{n^2+an+2}-\sqrt{n^2+bn+1}\right)=1$. Khẳng định nào sau đây đúng?
	\choice 
		 {$a+b=2$}
		 {\True $a-b=2$}
		{$a+b=1$} 
		 {$a-b=1$}
	\loigiai{ 
		Ta có \allowdisplaybreaks\begin{eqnarray*}
		 &&\lim\left(\sqrt{n^2+an+2}-\sqrt{n^2+bn+1}\right)=1\\
		&\Leftrightarrow&\lim\dfrac{\left(n^2+an+2\right)-\left(n^2+bn+1\right)}{\sqrt{n^2+an+1}+\sqrt{n^2+bn+2}}=1\\ 
		&\Leftrightarrow  &\lim\dfrac{(a-b)n+1}{\sqrt{n^2+an+2}+\sqrt{n^2+bn+1}}=1\\
		 &\Leftrightarrow &\lim\dfrac{(a-b)+\dfrac{1}{n}}{\sqrt{1+\dfrac{a}{n}+\dfrac{2}{n^2}}+\sqrt{1+\dfrac{b}{n}+\dfrac{1}{n^2}}}=1\\
		 &\Leftrightarrow &\dfrac{a-b}{2}=1\\
		&\Leftrightarrow &a-b=2. 
	\end{eqnarray*}
} 
\end{ex}

% nguồn: Ok_Toan11_TracNghiem_Lan2
\begin{ex}%[1D3V1-3]%[ma:2010019d]
	Có bao nhiêu giá trị nguyên của $a$ để $\lim\left(\sqrt{n^2-4n+7}+a-n\right)=0$? 
	 \choice
		 {$3$}
		{\True $1$} 
		{$2$} 
		 {$0$}
	 \loigiai{
		Ta có \allowdisplaybreaks\begin{eqnarray*}
		\lim\left(\sqrt{n^2-4n+7}+a-n\right)&=&\lim\dfrac{-4n+7+2an-a^2}{\sqrt{n^2-4n+7}-(a-n)}\\ 
		 &=&\lim\dfrac{2a-4+\dfrac{7-a^2}{n}}{\sqrt{1-\dfrac{4}{n}+\dfrac{7}{n^2}}-\dfrac{a}{n}+1}=a-2.
	\end{eqnarray*}
	 Để $\lim\left(\sqrt{n^2-4n+7}+a-n\right)=0$ thì $a-2=0\Leftrightarrow a=2$.\\
	Vậy có một giá trị nguyên  của $a$ thỏa mãn bài toán.
 }
\end{ex}

% nguồn: Ok_Toan11_TracNghiem_Lan2
\begin{ex}%[1D3C1-3]%[ma:143be1b4]
	Tính  $\lim n\left(\sqrt{4n^2+1}-\sqrt[3]{8n^3+n}\right)$.
	\choice 
		 {$+\infty$}
		{$1$} 
		 {$-\infty$}
		{\True  $\dfrac{1}{6}$}
	\loigiai{Ta có \allowdisplaybreaks\begin{eqnarray*}
		\lim n\left(\sqrt{4n^2+1}-\sqrt[3]{8n^3+n}\right)&=&\lim  n\left[\left(\sqrt{4n^2+1}-2n\right)+\left(2n-\sqrt[3]{8n^3+n}\right)\right]\\
		 &=&\lim\left[n\left(\sqrt{4n^2+1}-2n\right)+n\left(2n-\sqrt[3]{8n^3+n}\right)\right]\\
		&=&\lim  n\left(\sqrt{4n^2+1}-2n\right)+\lim n\left(2n-\sqrt[3]{8n^3+n}\right).
	\end{eqnarray*}
	Ta tính được \begin{itemize}
	\item  $\lim n\left(\sqrt{4n^2+1}-2n\right)=\lim\dfrac{n}{\sqrt{4n^2+1}+2n}=\lim\dfrac{1}{\sqrt{4+\dfrac{1}{n^2}}+2}=\dfrac{1}{4}$,
	\item $\begin{aligned}
	 \lim n\left(2n-\sqrt[3]{8n^3+n}\right)&=\lim\dfrac{-n^2}{4n^2+2n\sqrt[3]{8n^3+n}+\left(\sqrt[3]{8n^3+n}\right)^2}\\
	&=\lim\dfrac{-1}{4+2\sqrt[3]{8+\dfrac{1}{n^2}}+\left(\sqrt[3]{8+\dfrac{1}{n^2}}\right)^2}=-\dfrac{1}{12}. 
	\end{aligned}$
	\end{itemize}
	Vậy $\lim n\left(\sqrt{4n^2+1}-\sqrt[3]{8n^3+n}\right)=\dfrac{1}{4}-\dfrac{1}{12}=\dfrac{1}{6}$. 
} 
\end{ex}


%%%%%%%%%%%%%%%% ĐÚNG/SAI %%%%%%%%%%%%%%%%

% nguồn: Ok_dungsai_lop11_Lan1
\begin{ex}%[1D3N1-3]%[ma:1536bc0f]
 Mỗi kết quả giới hạn dưới đây đúng  hay sai?
  \choiceTF
  {\True $\lim\limits_{n \to +\infty}  \left(\sqrt{n}-\sqrt{n+1}\right)=0$}
  {$\lim\limits_{n \to +\infty} \left(\sqrt{n+2}-\sqrt{n+1}\right)=1$} 
  {\True $\lim\limits_{n \to +\infty}  \left(\sqrt{n^2-4n}-n\right)=-2$}
  {$\lim\limits_{n \to  +\infty} \left(2n-\sqrt{4n^2+n}\right)=0$}
 \loigiai{ 
  \begin{itemchoice}
   \itemch $\lim\limits_{n \to +\infty}  \left(\sqrt{n}-\sqrt{n+1}\right)
   =\lim\limits_{n  \to +\infty} \dfrac{-1}{\sqrt{n}+\sqrt{n+1}}
   =\lim\limits_{n \to +\infty} \dfrac{-1}{\sqrt{n}\left(1+\sqrt{1+\dfrac{1}{n}}\right)}=0$. 
   \itemch $\lim\limits_{n \to +\infty} \left(\sqrt{n+2}-\sqrt{n+1}\right)=\lim\limits_{n \to +\infty} \dfrac{1}{\sqrt{n+1}+\sqrt{n+1}} 
   =\lim\limits_{n \to +\infty}  \dfrac{1}{\sqrt{n}\left(\sqrt{1+\dfrac{2}{n}}+\sqrt{1+\dfrac{1}{n}}\right)}=0$.
   \itemch $\lim\limits_{n \to +\infty} \left(\sqrt{n^2-4n}-n\right) 
   =\lim\limits_{n \to +\infty} \dfrac{-4n}{\sqrt{n^2-4n}+n} 
   =\lim\limits_{n \to  +\infty} \dfrac{-4}{\sqrt{1-\dfrac{4}{n}}+1}
   =-2$. 
   \itemch $\lim\limits_{n \to +\infty}  \left(2n-\sqrt{4n^2+n}\right)
   =\lim\limits_{n \to  +\infty} \dfrac{-n}{2n+\sqrt{4n^2+n}}
   =\lim\limits_{n \to +\infty}  \dfrac{-1}{2+\sqrt{4+\dfrac{1}{n^2}}}=-\dfrac{1}{4}$.
  \end{itemchoice}
 } 
\end{ex}

% nguồn: Ok_dungsai_lop11_Lan1
\begin{ex}%[1D3H1-3]%[ma:24371bd3]
 Mỗi kết  quả giới hạn dưới đây đúng hay sai?
 \choiceTF 
   {$\lim\limits_{n \to +\infty} \left(\sqrt{n+3}-\sqrt{n+1}\right)=2$}
  {$\lim\limits_{n \to  +\infty} \left(n-\sqrt{n^2-1}\right)=-1$}
  {\True  $\lim\limits_{n \to +\infty} \dfrac{\sqrt{n+3}-\sqrt{n}}{n+2}=0$}
  {$\lim\limits_{n  \to +\infty}\left(\sqrt{n^2+2}-n\right)n=0$}
 \loigiai{ 
  \begin{itemchoice}
   \itemch $\lim\limits_{n \to  +\infty} \left(\sqrt{n+3}-\sqrt{n+1}\right)
   =\lim\limits_{n  \to +\infty} \dfrac{2}{\sqrt{n+3}+\sqrt{n+1}}
    =\lim\limits_{n \to +\infty} \dfrac{1}{\sqrt{n}\left(\sqrt{1+\dfrac{3}{n}}+\sqrt{1+\dfrac{1}{n}}\right)}=0$.
   \itemch $\lim\limits_{n \to +\infty} \left(n-\sqrt{n^2+1}\right)=\lim\limits_{n \to  +\infty} \dfrac{1}{n+\sqrt{n^2-1}}=\lim\limits_{n \to \infty} \dfrac{1}{n\left(1+\sqrt{1-\dfrac{1}{n^2}}\right)}=0$.
   \itemch $\lim\limits_{n \to +\infty}  \dfrac{\sqrt{n+3}-\sqrt{n}}{n+2}
   =\lim\limits_{n \to +\infty} \dfrac{3}{\left(\sqrt{n+3}+\sqrt{n}\right)(n+2)} 
   =\lim\limits_{n  \to +\infty} \dfrac{\dfrac{3}{n\sqrt{n}}}{\left(\sqrt{1+\dfrac{3}{n}}+\sqrt{n}\right)(1+\dfrac{2}{n})}=0$.
    \itemch $\lim\limits_{n \to +\infty}\left(\sqrt{n^2+2}-n\right)n
   =\lim\limits_{n  \to +\infty} \dfrac{2n}{\sqrt{n^2+2}+n}
   =\lim\limits_{n \to +\infty}  \dfrac{2}{\sqrt{1+\dfrac{2}{n^2}}+1}=1
    $.
  \end{itemchoice}
  }
\end{ex}

% nguồn: Ok_dungsai_lop11_Lan2
\begin{ex}%[1D3H1-3]%[ma:c60effec]
	 Đặt $I=\lim \left(\sqrt{n^2+a^2n}-\sqrt{n^2+(a+2) n+1}\right)$.
	 \choiceTF
		{Ta  biến đổi được $I=\lim \dfrac{n-1}{\sqrt{n^2+n}+\sqrt{n^2+1}}$}
		{Nếu $I=0$ thì có $3$ giá trị $a$ thỏa mãn} 
		{Nếu  $I=0$ thì tổng các giá trị $a$ tìm được bằng $1$}
		{Có  $2$ giá trị $a$ nguyên để $I=1$}
	\loigiai{ 
		\begin{itemchoice}
			\itemch 
			Ta có 
			\begin{eqnarray*}
				\lim \left(\sqrt{n^2+a^2n}-\sqrt{n^2+(a+2) n+1}\right)&=&\lim \dfrac{\left(a^2-a-2\right)  n-1}{\sqrt{n^2+a^2n}+\sqrt{n^2+(a+2)n+1}}\\
				&=&\lim  \dfrac{a^2-a-2-\dfrac{1}{n}}{\sqrt{1+\dfrac{a^2}{n}}+\sqrt{1+\dfrac{1}{n^2}+\dfrac{a+2}{n}}}\\
				 &=&\dfrac{a^2-a-2}{2}.
			\end{eqnarray*}
			\itemch Khi $I=0$ thì  $\dfrac{a^2-a-2}{2}=0$ hay $a^2-a-2=0\Leftrightarrow\hoac{&a=-1\\&a=2.}$
			\itemch Nếu  $I=0$ thì tổng các giá trị $a$ tìm được bằng $1$.
			 \itemch
			 Khi $I=1$ thì $\dfrac{a^2-a-2}{2}=1$ hay $a^2-a-4=0\Leftrightarrow\hoac{&a=\dfrac{1-\sqrt{17}}{2} \\&a=\dfrac{1+\sqrt{17}}{2}.}$
		\end{itemchoice}
	 }
\end{ex}

% nguồn: Ok_dungsai_lop11_Lan2
\begin{ex}%[1D3H1-3]%[ma:8c898fcc]
	Cho dãy số $(u_n)$ với $u_n=\dfrac{4n^2+n+2}{a n^2+5}$, với $a\ne  0$.
	 \choiceTF
		{\True  Để dãy số đã cho có giới hạn bằng $2$, giá trị của $a=2$}
		{Để dãy số đã cho có giới hạn bằng  $1$, giá trị của $a=3$}
		{Để dãy số  đã cho có giới hạn bằng $3$, giá trị của $a$ là một số nguyên}
		{\True Để dãy số đã cho  có giới hạn bằng $-2$, giá trị của $a=-2$}
	 \loigiai{
		\begin{itemchoice}
			\itemch Ta có 
			\begin{eqnarray*}
				\lim \dfrac{4n^2+n+2}{a  n^2+5}&=&\lim \dfrac{4+\dfrac{1}{n}+\dfrac{2}{n^2}}{a+\dfrac{5}{n^2}}\\
				 &=&\dfrac{4}{a}.
			\end{eqnarray*}
			 Khi đó $2=\dfrac{4}{a} \Leftrightarrow a=2$.
			\itemch Khi đó $1=\dfrac{4}{a}  \Leftrightarrow a=4$.\\
			\itemch Khi đó  $3=\dfrac{4}{a} \Leftrightarrow a=\dfrac{4}{3}$.\\
			\itemch Khi đó $-2=\dfrac{4}{a} \Leftrightarrow a=-2$. 
		\end{itemchoice}
	 }
\end{ex}

% nguồn: Ok_dungsai_lop11_Lan2
\begin{ex}%[1D3V1-3]%[ma:59922f40]
	Cho giới  hạn $L=\lim \left(\sqrt{n^2+a^2n}-\sqrt{n^2+\left(a+2\right)n+1}\right)\cdot$
	\choiceTF 
		{\True  $a=0$ thì $L=-1$}
		{$a=-2$ thì $L=3$} 
		 {\True $a=2$ thì $L=0$}
		{\True Có $2$ giá trị của  $a$ để $\lim \left(\sqrt{n^2+a^2n}-\sqrt{n^2+(a+2)n+1}\right)=0$}
	\loigiai{ 
		\begin{itemchoice}
			 \itemch Khi $a=0$, ta có $L=\lim \left(n-\sqrt{n^2+2n+1}\right)
			=\lim  \left[n-(n+1)\right]=\lim (-1)=-1$.
			\itemch Khi $a=-2$, ta  có\\
			$\begin{aligned}
			L=\lim  \left(\sqrt{n^2+4n}-\sqrt{n^2+1}\right)
			=&\lim \dfrac{4n-1}{\sqrt{n^2+4n}+\sqrt{n^2+1}} \\ 
			=&\lim  \dfrac{4-\dfrac{1}{n}}{\sqrt{1+\dfrac{4}{n}}+\sqrt{1+\dfrac{1}{n^2}}}=2.
		\end{aligned}$
		\itemch Khi $a=2$,  ta có\\
		$\begin{aligned}
		L=\lim  \left(\sqrt{n^2+4n}-\sqrt{n^2+4n+1}\right)
		=&\lim \dfrac{-1}{\sqrt{n^2+4n}+\sqrt{n^2+4n+1}} \\ 
		 =&\lim \dfrac{-\dfrac{1}{n}}{\sqrt{1+\dfrac{4}{n}}+\sqrt{1+\dfrac{4}{n}+\dfrac{1}{n^2}}}
		 =0.
	\end{aligned}$
	 \itemch Ta có\\
	$\begin{aligned}
	 L=&\lim \left(\sqrt{n^2+a^2n}-\sqrt{n^2+(a+2)n+1}\right)
	 =\lim \dfrac{(a^2-a-2)n-1}{\sqrt{n^2+a^2n}+\sqrt{n^2+(a+2)n+1}} \\
	 =&\lim \dfrac{(a^2-a-2)-\dfrac{1}{n}}{\sqrt{1+\dfrac{a^2}{n}}+\sqrt{1+\dfrac{(a+2)}{n}+\dfrac{1}{n^2}}}
	=\dfrac{a^2-a-2}{2}\cdot 
	\end{aligned}$ \\
	 Khi $L=0\Leftrightarrow \dfrac{a^2-a-2}{2}=0\Leftrightarrow a^2-a-2=0\Leftrightarrow \hoac{&a=-1 \\&a=2.}$ \\
	Suy ra có $2$ giá  trị của $a$ thỏa mãn.
	\end{itemchoice}
 }
\end{ex}


%%%%%%%%%%%%%%%% TRẢ LỜI NGẮN %%%%%%%%%%%%%%%%

% nguồn: Ok_TraLoiNgan_Lop11_Lan1
\begin{ex}%[1D3H1-3]%[ma:a81eb483]
	Tìm giới hạn sau: ${\lim  \left(\sqrt{3n^2+n+2}-\sqrt{3n^2+n+1}\right)}$.
	\shortans{0} 
	\loigiai{ 
		${\begin{aligned} & \lim \left(\sqrt{3n^2+n+2}-\sqrt{3n^2+n+1}\right)=\lim \dfrac{3n^2+n+2-\left(3n^2+n+1\right)}{\sqrt{3n^2+n+2}+\sqrt{3n^2+n+1}} \\ &=\lim \dfrac{1}{\sqrt{3n^2+n+2}+\sqrt{3n^2+n+1}}=0\left(\text{do } \lim \left(\sqrt{3n^2+n+2}+\sqrt{3n^2+n+1}\right)=+\infty\right).
	\end{aligned}}$
 }
\end{ex}

% nguồn: Ok_TraLoiNgan_Lop11_Lan1
\begin{ex}%[1D3H1-3]%[ma:815f49b0]
	Tìm giới hạn  sau: ${\lim \left(\sqrt{n^2+2n-2}-n\right)}$.
	 \shortans{1}
	\loigiai{ 
		$$\begin{aligned}
		 &\lim \left(\sqrt{n^2+2n-2}-n\right)=\lim \dfrac{n^2+2n-2-n^2}{\sqrt{n^2+2n-2}+n}=\lim \dfrac{2n-2}{n \sqrt{1+\dfrac{2}{n}-\dfrac{2}{n^2}}+n}\\
		&=\lim  \dfrac{n\left(2-\dfrac{2}{n}\right)}{n\left(\sqrt{1+\dfrac{2}{n}-\dfrac{2}{n^2}}+1\right)}=\lim \dfrac{2-\dfrac{2}{n}}{\sqrt{1+\dfrac{2}{n}-\dfrac{2}{n^2}}+1}=\dfrac{2}{\sqrt{1}+1}=1.
	\end{aligned}
	$$ 
 }
\end{ex}

% nguồn: Ok_TraLoiNgan_Lop11_Lan1
\begin{ex}%[1D3H1-3]%[ma:09ae7626]
	Tìm giới hạn sau: ${\lim  6\left(\sqrt[3]{n^3-2n^2}-n\right)}$.
	 \shortans{-4}
	\loigiai{ 
		$${\begin{aligned}
		& \lim \left(\sqrt[3]{n^3-2 n^2}-n\right)=\lim \dfrac{n^3-2 n^2-n^3}{\sqrt[3]{\left(n^3-2 n^2\right)^2}+n \sqrt[3]{n^3-2 n^2}+n^2} \\ 
		&=\lim \dfrac{-2 n^2}{\sqrt[3]{\left(n^3-2 n^2\right)^2}+n \sqrt[3]{n^3-2 n^2}+n^2}=\lim  \dfrac{-2}{\sqrt[3]{\left(1-\dfrac{2}{n}\right)^2}+\sqrt[3]{1-\dfrac{2}{n}}+1}=-\dfrac{2}{3}
	\end{aligned}
	 }$$
	Do đó $\lim 6\left(\sqrt[3]{n^3-2n^2}-n\right)=-4$. 
} 
\end{ex}

% nguồn: Ok_TraLoiNgan_Lop11_Lan2
\begin{ex}%[1D3H1-3]%[ma:35ea99eb]
	Tính giới hạn $L=\lim  \left(\sqrt{n^{24}+6n^{12}+1}-\sqrt[3]{n^{36}+3n^{24}+2}\right)$.
	 \shortans{2}
	 \loigiai{
		Ta có 
		\begin{eqnarray*}
			L=\lim  \left(\sqrt{n^{24}+6n^{12}+1}-\sqrt[3]{n^{36}+3n^{24}+2}\right)=L_1+L_2.
		\end{eqnarray*}
		Trong đó 
		\begin{eqnarray*}
			L_1&=&\lim \left(\sqrt{n^{24}+6 n^{12}+1}-n^{12}\right)\\ 
			&=&\lim \dfrac{6  n^{12}+1}{\sqrt{n^{24}+6 n^{12}+1}+n^{12}}\\
			 &=&\lim \dfrac{6+\dfrac{1}{n^{12}}}{\sqrt{1+\dfrac{6}{n^{12}}+\dfrac{1}{n^{24}}}+1}\\
			 &=&3.
		\end{eqnarray*}
		 Và
		\begin{eqnarray*}
			L_2&=&\lim  \left(n^{12}-\sqrt[3]{n^{36}+3 n^{24}+2}\right)\\
			&=&\lim \dfrac{-3  n^{24}-2}{n^{24}+n^{12} \sqrt[3]{n^{36}+3 n^{24}+2}+\left(\sqrt[3]{n^{36}+3 n^{24}+2}\right)^2}\\
			&=&\lim \dfrac{-3-\dfrac{2}{n^{24}}}{1+\sqrt[3]{1+\dfrac{3}{n^{12}}+\dfrac{2}{n^{36}}}+\left(\sqrt[3]{1+\dfrac{3}{n^{12}}+\dfrac{2}{n^{36}}}\right)^2}\\ 
			&=&-1. 
		\end{eqnarray*}
		 Vậy $L=3+(-1)=2$.
	} 
\end{ex}

% nguồn: Ok_TraLoiNgan_Lop11_Lan2
\begin{ex}%[1D3H1-3]%[ma:cfbcd6b4]
	Biết rằng giới hạn $\lim \left(1+3n-\sqrt{9n^2-n+7}\right)=\lim \left(a+\dfrac{n-7}{b. n+c. \sqrt{9n^2-n+7}}\right)$ với $a$; $b$;  $c$ là các số tự nhiên. Tính giá trị $|a-b|+c$.
	 \shortans{12}
	 \loigiai{
		Ta có 
		\begin{eqnarray*}
			\lim \left(1+3n-\sqrt{9n^2-n+7}\right)&=&\lim  \left(1+\dfrac{9n^2-9n^2+n-7}{3n+\sqrt{9n^2-n+7}}\right)\\
			 &=&\lim \left(1+\dfrac{n-7}{3n+\sqrt{9n^2-n+7}}\right).
		\end{eqnarray*}
		 Khi đó $\heva{&a=1\\&b=3\\&c=1} \Rightarrow|a-b|+c=3$.
	 }
\end{ex}

% nguồn: Ok_TraLoiNgan_Lop11_Lan2
\begin{ex}%[1D3H1-3]%[ma:257d214f]
	Biết giới  hạn $\lim \left(\sqrt[3]{n^3+3}-\sqrt{n^2+2}\right)=\lim \dfrac{a}{\sqrt[3]{\left(n^3+3\right)^2}+n^2+n \cdot \sqrt[3]{n^3+3}}-\lim \dfrac{b}{n+\sqrt{n^2+2}}$ với $a$; $b$ là các số nguyên dương. Tính $T=2a+3b$.
	 \shortans{12}
	\loigiai{ 
		Ta có 
		\begin{eqnarray*}
			\lim \left(\sqrt[3]{n^3+3}-\sqrt{n^2+2}\right)&=&\lim \left(\sqrt[3]{n^3+3}-n\right)+\lim  \left(n-\sqrt{n^2+2}\right)\\
			&=&\lim \dfrac{n^3+3-n^3}{\left(n^3+3\right)^{\tfrac{2}{3}}+n^2+n \cdot \sqrt[3]{n^3+3}}+\lim \dfrac{n^2-n^2-2}{n+\sqrt{n^2+2}}\\ 
			&=&\lim \dfrac{3}{\left(n^3+3\right)^{\tfrac{2}{3}}+n^2+n \cdot  \sqrt[3]{n^3+3}}-\lim \dfrac{2}{n+\sqrt{n^2+2}}\\
			&=&\lim \dfrac{3}{\sqrt[3]{\left(n^3+3\right)^2}+n^2+n  \cdot \sqrt[3]{n^3+3}}-\lim \dfrac{2}{n+\sqrt{n^2+2}}.
		\end{eqnarray*}
		Do  đó, $\heva{&a=3\\&b=2}\Rightarrow T=2a+3b=12$.
	} 
\end{ex}

% nguồn: Ok_TraLoiNgan_Lop11_Lan2
\begin{ex}%[1D3H1-3]%[ma:d4e846f6]
	Giới hạn dãy số $u_n=\dfrac{\sqrt{n^2+2n}-\sqrt{n^2+n}}{n}$ có dạng $\lim  \dfrac{\dfrac{a}{n}}{\sqrt{1+\dfrac{b}{n}}+\sqrt{1+\dfrac{c}{n}}}$ với $a$; $b$; $c$ là các số tự nhiên. Tính giá trị $a^2+b^2+c^2$.
	\shortans{6} 
	 \loigiai{
		Ta có 
		\begin{eqnarray*}
			\lim  u_n=\lim \dfrac{\sqrt{n^2+2n}-\sqrt{n^2+n}}{n}&=&\lim \dfrac{\left(n^2+2n\right)-\left(n^2+n\right)}{n\left(\sqrt{n^2+2n}+\sqrt{n^2+n}\right)}\\
			 &=&\lim \dfrac{n}{n\left(\sqrt{n^2+2n}+\sqrt{n^2+n}\right)}\\
			&=&\lim \dfrac{1}{\sqrt{n^2+2n}+\sqrt{n^2+n}}\\ 
			 &=&\lim \dfrac{\dfrac{1}{n}}{\sqrt{1+\dfrac{2}{n}}+\sqrt{1+\dfrac{1}{n}}}.
		\end{eqnarray*}
		Do  đó, $\heva{&a=1\\&b=2\\&c=1} \Rightarrow a^2+b^2+c^2=6$.
	 }
\end{ex}

% nguồn: Ok_TraLoiNgan_Lop11_Lan2
\begin{ex}%[1D3H1-3]%[ma:09dea702]
	Biết rằng giới hạn $\lim  \left(\sqrt{n^2-4n+5}+3-n\right)=\lim \left(\dfrac{5-a \cdot n}{\sqrt{n^2-4n+5}+n}+b\right)$ với $a$; $b$ là các số tự nhiên. Tính giá trị $S=a^b$.
	\shortans{64} 
	 \loigiai{
		 Ta có
		\begin{eqnarray*}
			 \lim \left(\sqrt{n^2-4n+5}+3-n\right)&=&\lim \left(\dfrac{\left(\sqrt{n^2-4n+5}-n\right) \cdot\left(\sqrt{n^2-4n+5}+n\right)}{\sqrt{n^2-4n+5}+n}+3\right)\\
			&=&\lim \left(\dfrac{n^2-4n+5-n^2}{\sqrt{n^2-4n+5}+n}+3\right)\\ 
			&=&\lim \left(\dfrac{-4n+5}{\sqrt{n^2-4n+5}+n}+3\right). 
		\end{eqnarray*}
		 Khi đó $\heva{&a=4\\&b=3} \Rightarrow S=a^b=64$.
	} 
\end{ex}

% nguồn: Ok_TraLoiNgan_Lop11_Lan1
\begin{ex}%[1D3V1-3]%[ma:b02629a5]
	 Tính $\lim\left(\sqrt{n^2+2n}-n\right)$.
	\shortans{1} 
	\loigiai{ 
		\allowdisplaybreaks 
		\begin{eqnarray*}
			 \lim\left(\sqrt{n^2+2n}-n\right)&=&\lim\dfrac{\left(\sqrt{n^2+2n}-n\right)\left(\sqrt{n^2+2n}+n\right)}{\left(\sqrt{n^2+2n}+n\right)}\\
			&=&\lim  \dfrac{n^2+2n-n^2}{\left(\sqrt{n^2+2n}+n\right)}\\
			 &=&\lim \dfrac{2n}{\left(\sqrt{n^2+2n}+n\right)}\\
			&=&\lim \dfrac{2}{\sqrt{1+\dfrac{2}{n}}+1}\\ 
			&=&1. 
		\end{eqnarray*}
	} 
\end{ex}

% nguồn: Ok_TraLoiNgan_Lop11_Lan2
\begin{ex}%[1D3V1-3]%[ma:30eb0c7d]
	Biết $\lim \dfrac{\sqrt{16n^4-4}-\sqrt[3]{n^6+\pi}}{an^2+\pi}=-12$ với $a\ne 0$. Biết $a=-\dfrac{m}{n}$ tối giản  và $m, n\in \mathbb{N}^*$. Tính $m+n$.
	\shortans{5} 
	 \loigiai{
		Ta có $\lim \dfrac{\sqrt{16n^4-4}-\sqrt[3]{n^6+\pi}}{an^2+\pi}=\lim \dfrac{\sqrt{16-\dfrac{4}{n^4}}-\sqrt[3]{1+\dfrac{\pi}{n^6}}}{a+\dfrac{\pi}{n^2}}=\dfrac{3}{a}$.\\ 
		Suy ra $\dfrac{3}{a}=-12\Rightarrow  a=-\dfrac{1}{4}\cdot$
	} 
\end{ex}

% nguồn: Ok_TraLoiNgan_Lop11_Lan2
\begin{ex}%[1D3V1-3]%[ma:b4b4c964]
	Giả sử $\lim \dfrac{\sqrt{9n^2-4}-\sqrt[3]{n^3+n^2}}{5n+11}=\dfrac{a}{b}$ với $a$, $b\in \mathbb{N}$ và $\dfrac{a}{b}$ tối giản.  Giá trị của $T=2a-3b$ bằng
	 \shortans{-11}
	 \loigiai{
		Ta  có $\lim \dfrac{\sqrt{9n^2-4}-\sqrt[3]{n^3+n^2}}{5n+11}
		=\lim  \dfrac{\sqrt{9-\dfrac{4}{n^2}}-\sqrt[3]{1+\dfrac{1}{n}}}{5+\dfrac{11}{n}}=\dfrac{2}{5}$.\\
		Suy ra $a=2$, $b=5$. Do đó  $T=2\cdot 2-3\cdot 5=-11$.
	 }
\end{ex}

% nguồn: Ok_TraLoiNgan_Lop11_Lan2
\begin{ex}%[1D3V1-3]%[ma:490ea9d8]
	 Tính $a+b$ biết
	 \[
	\lim \left( an - \sqrt{n^2 + bn +  2} \right) = 1.
	\] 
	\shortans{-1} 
	\loigiai{ 
		Nếu $a \neq  1$ thì giới hạn là vô hạn, nên $a = 1$. \\
		Tiếp  tục
		 \[
		\lim \left(n - \sqrt{n^2 + bn + 2}\right)  = -\dfrac{b}{2} = 1 \Rightarrow b = -2.
		 \]
		Vậy $a + b  = -1$.
	 }
\end{ex}

% nguồn: Ok_TraLoiNgan_Lop11_Lan2
\begin{ex}%[1D3V1-3]%[ma:ea7a754b]
	Biết $\lim \dfrac{\sqrt{n^4-4}-\sqrt[3]{27n^6+5}+9n+1}{an^2+n+1}=-1$ với  $a\ne 0$. Tìm $a$.
	\shortans{2} 
	 \loigiai{
		Ta có  $\lim \dfrac{\sqrt{n^4-4}-\sqrt[3]{27n^6+5}+9n+1}{an^2+n+1}
		 =\lim \dfrac{\sqrt{1-\dfrac{4}{n^4}}-\sqrt[3]{27+\dfrac{5}{n^6}} +\dfrac{9}{n}+\dfrac{1}{n^2}}{a+\dfrac{1}{n}+\dfrac{1}{n^2}}=\dfrac{-2}{a}$.\\
		Suy ra $\dfrac{-2}{a}=-1\Rightarrow a=2$. 
	 }
\end{ex}


%%%%%%%%%%%%%%%% TỰ LUẬN %%%%%%%%%%%%%%%%

% nguồn: Ok_TuanLuan_Toan11_Lan2
\begin{ex}%[1D3H1-3]%[ma:3051e33b]
	Cho dãy số $\left(u_n \right)$ với $u_n=\dfrac{1}{1\sqrt{2}+2\sqrt{1}}+\dfrac{1}{2\sqrt{3}+3\sqrt{2}}+\cdots+\dfrac{1}{n\sqrt{n+1}+(n+1)\sqrt{n}}$. Tính $\lim \left(u_n-1 \right)$. 
	 \loigiai{
		 Ta có
		\begin{eqnarray*}
			u_n&=&\dfrac{1}{n\sqrt{n+1}+(n+1)\sqrt{n}}\\ 
			&=&\dfrac{1}{\sqrt{n}\sqrt{n+1} \left(\sqrt{n}+\sqrt{n+1} \right)}\\ 
			&=&\dfrac{\sqrt{n+1}-\sqrt{n}}{\sqrt{n}\sqrt{n+1}}=\dfrac{1}{\sqrt{n}}-\dfrac{1}{\sqrt{n+1}}\\ 
			 &=&\dfrac{1}{1\sqrt{2}+2\sqrt{1}}+\dfrac{1}{2\sqrt{3}+3\sqrt{2}}+ \cdots +\dfrac{1}{n\sqrt{n+1}+(n+1)\sqrt{n}}\\
			 &=&\dfrac{1}{\sqrt{1}}-\dfrac{1}{\sqrt{2}}+\dfrac{1}{\sqrt{2}}-\dfrac{1}{\sqrt{3}}+\cdots +\dfrac{1}{\sqrt{n}}-\dfrac{1}{\sqrt{n+1}}\\
			 &=&1-\dfrac{1}{\sqrt{n+1}}.
		\end{eqnarray*}
		Vậy $\lim \left(u_n-1 \right)=\lim \left(-\dfrac{1}{\sqrt{n+1}} \right)=0$. 
	} 
\end{ex}

% nguồn: Ok_TuanLuan_Toan11_Lan2
\begin{ex}%[1D3H1-3]%[ma:c136edd8]
	Cho dãy số $\left(u_n \right)$  với $u_n=\sqrt[3]{n+2}-\sqrt[3]{n}$. Tính $\lim u_n$.
	 \loigiai{
		Ta  có
		\begin{eqnarray*}
			u_n&=&\sqrt[3]{n+2}-\sqrt[3]{n}\\ 
			&=&\dfrac{n+2-n}{\left(\sqrt[3]{n+2} \right)^2+\sqrt[3]{n+2} \cdot \sqrt[3]{n}+\left(\sqrt[3]{n}  \right)^2}\\
			&=&\dfrac{2}{\left(\sqrt[3]{n\left(1+\dfrac{2}{n} \right)} \right)^2+\sqrt[3]{n\left(1+\dfrac{2}{n} \right)}  \cdot \sqrt[3]{n}+\left(\sqrt[3]{n} \right)^2}\\
			 &=& \dfrac{2}{\left(\sqrt[3]{n} \right)^2 \left[\left(\sqrt[3]{1+\dfrac{2}{n}} \right)^2+\sqrt[3]{1+\dfrac{2}{n}}+1 \right]}.
		\end{eqnarray*}
		Vì $\lim \dfrac{2}{3\left(\sqrt[3]{n}  \right)^2}=0$ và $\lim \dfrac{1}{\left(\sqrt[3]{1+\dfrac{2}{n}} \right)^2+\sqrt[3]{1+\dfrac{2}{n}}+1}=\dfrac{1}{3}$.\\
	Do đó $\lim  u_n=0$.}
\end{ex}

% nguồn: Ok_TuanLuan_Toan11_Lan2
\begin{ex}%[1D3H1-3]%[ma:b1cc157a]
	Cho  dãy số $\left(u_n \right)$ với $u_n=\dfrac{\sqrt{4n^2+1}-2n}{\sqrt{n^2+4n+1}-n}$. Tính $\lim u_n$.
	\loigiai{ 
		Xử  lí tử số \[\sqrt{4n^2+1}-2n=\dfrac{4n^2+1-4n^2}{\sqrt{4n^2+1}+2n}=\dfrac{1}{\sqrt{4n^2+1}+2n}.\] Xử lí mẫu số \[\dfrac{1}{\sqrt{n^2+4n+1}-n}=\dfrac{\sqrt{n^2+4n+1}+n}{n^2+4n+1-n^2}=\dfrac{\sqrt{n^2+4n+1}+n}{4n+1}.\]
		 Từ đó, ta có
		\begin{eqnarray*}
			\lim u_n&=&\lim \dfrac{\sqrt{n^2+4n+1}+n}{\left(\sqrt{4n^2+1}+2n  \right)(2n+1)}\\
			 &=&\lim \dfrac{\sqrt{n^2\left(1+\dfrac{4}{n}+\dfrac{1}n^2 \right)}+n}{\left(\sqrt{n^2\left(4+\dfrac{1}n^2 \right)}+2n \right)(2n+1)}\\
			&=&\lim \dfrac{n\left[\sqrt{1+\dfrac{4}{n}+\dfrac{1}n^2}+1 \right]}{n\left(\sqrt{4+\dfrac{1}n^2}+2 \right)n\left(2+\dfrac{1}{n}  \right)}\\
			&=&\lim \dfrac{\sqrt{1+\dfrac{4}{n}+\dfrac{1}n^2}+1}{n\left(\sqrt{4+\dfrac{1}n^2}+2 \right)\left(2+\dfrac{1}{n} \right)}\\ 
			 &=&\lim \dfrac{2}{n(2+2)2}=\lim \dfrac{1}{4n}=0.
		\end{eqnarray*}
	Do đó $\lim  u_n=0$.}
\end{ex}

% nguồn: Ok_TuanLuan_Toan11_Lan2
\begin{ex}%[1D3H1-3]%[ma:11b7c3a1]
	Cho dãy số $\left(u_n \right)$  với $u_n=\dfrac{1+2+3+4+\cdots+n}{\left(1+3+3^2+3^3+\cdots+3^n \right) \cdot (n+1)}$. Tính $\lim u_n$.
	 \loigiai{
		Xét tử số, ta thấy $1$, $2$, $3$, $4$, $\cdots$, $n$ là một dãy số thuộc cấp  số cộng có $n$ số hạng với $u_1=1$, $d=1$. Tổng $n$ số hạng của cấp số cộng \[S_n=\dfrac{\left(u_1+u_n \right)n}{2}=\dfrac{(1+n)n}{2}.\]
		Xét mẫu số, ta thấy $1$, $3$, $3^2$, $3^3$, $\cdots$, $3^n$ là một dãy số thuộc cấp số nhân có $(n+1)$ số hạng với $u_1=1$,  $q=3$.\\
		Tổng $(n+1)$  số hạng của cấp số nhân \[S_{n+1}=u_1 \cdot \dfrac{1-q^{n+1}}{1-q}=\dfrac{1-3^{n+1}}{1-3}=\dfrac{3^{n+1}-1}{2}\]
		\[u_n=\dfrac{n}{3^{n+1}-1}=\dfrac{n}{3 \cdot  3^n-1}.\]
		Bằng quy nạp ta luôn có $n < 2^n$, $\forall n\in \mathbb{N}^{*}$ và $3^n >  1$, $\forall n\in \mathbb{N}^{*}$\\
		Suy ra  $\left|u_n \right|=\left|\dfrac{n}{3 \cdot 3^n-1} \right| < \dfrac{n}{3^n} < \dfrac{2^n}{3^n}=\left(\dfrac{2}{3} \right)^n$.\\
	Vì $\lim  \left(\dfrac{2}{3} \right)^n=0$ nên $\lim u_n=0$. }
\end{ex}

% nguồn: Ok_TuanLuan_Toan11_Lan2
\begin{ex}%[1D3V1-3]%[ma:263ae6bf]
	Tìm $\lim  \left(\sqrt[3]{n^3+3n^2}-n\right)$.
	 \loigiai{
		 \allowdisplaybreaks
		\begin{eqnarray*}
			 \lim \left(\sqrt[3]{n^3+3n^2}-n\right)
			&=& \lim \dfrac{\left(\sqrt[3]{n^3+3n^2}-n\right)\left(\sqrt[3]{{\left(n^3+3n^2\right)}^2}+n\sqrt[3]{n^3+3n^2}+n^2\right)}{\sqrt[3]{{\left(n^3+3n^2\right)}^2}+n\sqrt[3]{n^3+3n^2}+n^2} \\ 
			&=& \lim \dfrac{3n^2}{\sqrt[3]{{\left(n^3+3n^2\right)}^2}+n\sqrt[3]{n^3+3n^2}+n^2} \\ 
			&=& \lim \dfrac{3}{\sqrt[3]{{\left(1+\dfrac{3}{n}\right)}^2}+\sqrt[3]{1+{\dfrac{3}{n}}^2}+1} \\ 
			&=& 1. 
		\end{eqnarray*}
	} 
\end{ex}

% nguồn: Ok_TuanLuan_Toan11_Lan2
\begin{ex}%[1D3V1-3]%[ma:461d4b13]
	Tìm  $\lim \dfrac{2n-\sqrt{4n^2+n}}{n+\sqrt[3]{4n^2-n^3}}$.
	\loigiai{ 
		Ta  có
		 \allowdisplaybreaks
		\begin{eqnarray*}
			 \lim \left(2n-\sqrt{4n^2+n}\right)
			&=& \lim \dfrac{\left(2n-\sqrt{4n^2+n}\right)\left(2n+\sqrt{4n^2+n}\right)}{2n+\sqrt{4n^2+n}} \\ 
			&=&  \lim \dfrac{-n}{2n+\sqrt{4n^2+n}} \\
			&=& \lim \dfrac{-1}{2+\sqrt{4+\dfrac{1}{n}}}  \\
			&=& -\dfrac{1}{4}. 
		\end{eqnarray*}
		 Ta có
		\begin{eqnarray*}
			\lim \left(n+\sqrt[3]{4n^2-n^3}\right) 
			 &=& \lim \dfrac{\left(n+\sqrt[3]{4n^2-n^3}\right)\left(n^2-n\sqrt[3]{4n^2-n^3}+\sqrt[3]{\left(4n^2-n^3\right)^2}\right)}{n^2-n\sqrt[3]{4n^2-n^3}+\sqrt[3]{\left(4n^2-n^3\right)^2}} \\
			&=& \lim  \dfrac{4n^2}{n^2-n\sqrt[3]{4n^2-n^3}+\sqrt[3]{\left(4n^2-n^3\right)^2}} \\
			&=&  \lim \dfrac{4}{1-\sqrt[3]{\dfrac{4}{n}-1}+\sqrt[3]{\left(\dfrac{4}{n}-1\right)^2}} \\
			&=& \dfrac{4}{1+1+1}  \\
			&=&  \dfrac{4}{3}.
		\end{eqnarray*}
		Suy ra $\lim  \dfrac{2n-\sqrt{4n^2+n}}{n+\sqrt[3]{4n^2-n^3}} = \left(-\dfrac{1}{4}\right) : \dfrac{4}{3} = -\dfrac{3}{16}$.
	} 
\end{ex}

% nguồn: Ok_TuanLuan_Toan11_Lan2
\begin{ex}%[1D3V1-3]%[ma:25814a5d]
	Tính  giới hạn $\lim\dfrac{3\sin n-4\cos n}{2n^2+1}$.
	\loigiai{ 
		 Vì $\left|3\sin n-4\cos n\right|\le\sqrt{\left(3^2+4^2\right)\left(\sin^2n+\cos^2 n\right)}=5$ (Bất đẳng thức Bunhia-Copski).\\
		Nên $ 0\le\left|\dfrac{3\sin  n-4\cos n}{2n^2+1}\right|\le\dfrac{5}{2n^2+1}$.\\
	Mà $\lim\dfrac{5}{2n^2+1}=\lim\dfrac{\dfrac{5}{n^2}}{2+\dfrac{1}{n^2}}=0$ nên $\lim\dfrac{3\sin  n-4\cos n}{2n^2+1}=0$.}
\end{ex}

% nguồn: Ok_TuanLuan_Toan11_Lan2
\begin{ex}%[1D3V1-3]%[ma:671cfb58]
	Tìm $\lim  \dfrac{\sqrt{n^2+n}-n}{\sqrt{4n^2+3n}-2n}$.
	\loigiai{ 
		 Ta có
		\allowdisplaybreaks 
		\begin{eqnarray*}
			\lim  \left(\sqrt{n^2+n}-n\right)
			 &=& \lim \dfrac{\left(\sqrt{n^2+n}-n\right)\left(\sqrt{n^2+n}+n\right)}{\sqrt{n^2+n}+n} \\
			&=&  \lim \dfrac{n}{\sqrt{n^2+n}+n} \\
			&=&  \lim \dfrac{1}{\sqrt{1+\dfrac{1}{n}}+1} \\
			&=& \dfrac{1}{2}. 
		\end{eqnarray*}
		 Ta có
		\begin{eqnarray*}
			\lim  \left(\sqrt{4n^2+3n}-2n\right)
			&=& \lim \dfrac{\left(\sqrt{4n^2+3n}-2n\right)\left(\sqrt{4n^2+3n}+2n\right)}{\sqrt{4n^2+3n}+2n}  \\
			&=& \lim  \dfrac{3n}{\sqrt{4n^2+3n}+2n} \\
			&=& \lim \dfrac{3}{\sqrt{4+\dfrac{3}{n}}+2} \\ 
			&=&  \dfrac{3}{4}.
		\end{eqnarray*}
		Suy ra $\lim \dfrac{\sqrt{n^2+n}-n}{\sqrt{4n^2+3n}-2n} = \dfrac{1}{2}  : \dfrac{3}{4} = \dfrac{2}{3}$.
	} 
\end{ex}

% nguồn: Ok_TuanLuan_Toan11_Lan2
\begin{ex}%[1D3V1-3]%[ma:27da499a]
	 Tìm $\lim \left(n+\sqrt[3]{4n^2-n^3}\right)$.
	\loigiai{ 
		 Ta có
		\allowdisplaybreaks 
		\begin{eqnarray*}
			\lim \left(n+\sqrt[3]{4n^2-n^3}\right) 
			&=& \lim \dfrac{\left(n+\sqrt[3]{4n^2-n^3}\right)\left(n^2-n\sqrt[3]{4n^2-n^3}+\sqrt[3]{{\left(4n^2-n^3\right)}^2}\right)}{n^2-n\sqrt[3]{4n^2-n^3}+\sqrt[3]{{\left(4n^2-n^3\right)}^2}} \\ 
			&=& \lim  \dfrac{4n^2}{n^2-n\sqrt[3]{4n^2-n^3}+\sqrt[3]{{\left(4n^2-n^3\right)}^2}} \\
			&=& \lim  \dfrac{4}{1-\sqrt[3]{\dfrac{4}{n}-1}+\sqrt[3]{{\left(\dfrac{4}{n}-1\right)}^2}} \\
			&=& \dfrac{4}{1-1+1}  \\
			&=& 4. 
		\end{eqnarray*}
	 }
\end{ex}

% nguồn: Ok_TuanLuan_Toan11_Lan2
\begin{ex}%[1D3V1-3]%[ma:1f5ee696]
	Tính giới hạn $\lim\left(1-\dfrac{1}{2^2}\right)\left(1-\dfrac{1}{3^2}\right)\cdots\left(1-\dfrac{1}{n^2}\right)$. 
	\loigiai{ 
		 Ta có
		\begin{eqnarray*}
			 & &\left(1-\dfrac{1}{2^2}\right)\left(1-\dfrac{1}{3^2}\right)\left(1-\dfrac{1}{4^2}\right)\cdots\left(1-\dfrac{1}{\left(n-1\right)^2}\right)\left(1-\dfrac{1}{n^2}\right)\\
			&=&\dfrac{1}{2}\cdot\dfrac{3}{2}\cdot\dfrac{2}{3}\cdot\dfrac{4}{3}\cdot\dfrac{3}{4}\cdot\dfrac{5}{4}\cdots\dfrac{n-2}{n-1}\cdot\dfrac{n}{n-1}\cdot\dfrac{n+1}{n}=\dfrac{n+1}{2n}. 
		\end{eqnarray*}
	Suy ra $\lim\left(1-\dfrac{1}{2^2}\right)\left(1-\dfrac{1}{3^2}\right)\cdots\left(1-\dfrac{1}{n^2}\right)=\lim\dfrac{n+1}{2n}=\dfrac{1}{2}$.} 
\end{ex}

% nguồn: Ok_TuanLuan_Toan11_Lan2
\begin{ex}%[1D3V1-3]%[ma:1d941917]
	 Tìm $\lim \left(\sqrt{n^2+3n+5}-n\right)$.
	 \loigiai{
		\begin{itemize}
			\item \textbf{Cách  1.} Ta có
			\begin{eqnarray*}
				 \lim \left(\sqrt{n^2+3n+5}-n\right)
				&=& \lim \dfrac{\left(\sqrt{n^2+3n+5}-n\right)\left(\sqrt{n^2+3n+5}+n\right)}{\sqrt{n^2+3n+5}+n}\\ 
				&=& \lim  \dfrac{3n+5}{\sqrt{n^2+3n+5}+n}\\
				&=& \lim \dfrac{3+\dfrac{5}{n}}{\sqrt{1+\dfrac{3}{n}+\dfrac{5}{n^2}}+1}\\ 
				&=& \dfrac{3}{2}. 
			\end{eqnarray*}
			\item Nhân với một lượng liên hợp, sau đó rút gọn và làm như cách 2 ở  trên.
		\end{itemize}
		\textbf{Nhận xét:}  Khi nào sử dụng nhân với lượng liên hợp?
		\begin{itemize}
			\item  Khi $u_n=\sqrt{n^2+3n+5}-n=n\left(\sqrt{1+\dfrac{3}{n}+\dfrac{5}{n^2}}-1\right)$.\\
			Trong đó, $\lim n=+\infty$, $\lim \left(\sqrt{1+\dfrac{3}{n}+\dfrac{5}{n^2}}-1\right)=0$, khi đó $\lim u_n$ có dạng $+\infty \cdot 0$ (đây là một  dạng vô định) và ta không thể tính giới hạn của $u_n$ theo hướng này.
			\item Vậy khi nào thì chọn cách nhân với một lượng liên hợp? Cụ thể với $u_n=\sqrt{n^2+3n+5}-n$ xét ở trên trong  căn ta chỉ quan tâm đến biểu thức có chứa $n^2$ là cao nhất, còn lại bỏ hết, khi đó ta có thể xem $u_n=\sqrt{n^2}-n=0$, khi có điều này thì ta sẽ tìm giới hạn theo hướng nhân với một lượng liên hợp.
			\item Một ví dụ sau cho thấy ta không cần nhân với một lượng liên hợp. Ví dụ $u_n=\sqrt{2n^2+3n+5}-n$ xét ở trên trong căn ta chỉ quan tâm đến biểu thức có chứa $n^2$ là cao nhất, còn lại bỏ hết, khi đó ta có thể xem $u_n=\sqrt{2n^2}-n=n\left(\sqrt{2}-1\right)$, trong đó $\sqrt{2}-1>0$ và $\lim n=+\infty$,  nên giới hạn của $u_n$ là $+\infty$.
		\end{itemize}
		Cụ thể ta làm như sau: $\lim \left(\sqrt{2n^2+3n+5}-n\right) = \lim \left[n\left(\sqrt{2+\dfrac{3}{n}+\dfrac{5}{n^2}}-1\right)\right]  = +\infty$
	} 
\end{ex}

% nguồn: Ok_TuanLuan_Toan11_Lan2
\begin{ex}%[1D3V1-3]%[ma:937a809b]
	Tính giới hạn  $\lim\left(\dfrac{1}{\sqrt{4n^2+1}}+\dfrac{1}{\sqrt{4n^2+2}}+\cdots+\dfrac{1}{\sqrt{4n^2+n}}\right)$.
	\loigiai{ 
		Ta  có
		\begin{eqnarray*}
			 & &\dfrac{1}{\sqrt{4n^2}}+\dfrac{1}{\sqrt{4n^2}}+\cdots+\dfrac{1}{\sqrt{4n^2}}\le\dfrac{1}{\sqrt{4n^2+1}}+\dfrac{1}{\sqrt{4n^2+2}}+\cdots+\dfrac{1}{\sqrt{4n^2+n}}\\& \le &\dfrac{1}{\sqrt{4n^2+n}}+\dfrac{1}{\sqrt{4n^2+n}}+\cdots+\dfrac{1}{\sqrt{4n^2+n}}\\
			&\Leftrightarrow &\dfrac{n}{\sqrt{4n^2}}\le\dfrac{1}{\sqrt{4n^2+1}}+\dfrac{1}{\sqrt{4n^2+2}}+\cdots+\dfrac{1}{\sqrt{4n^2+n}}\le\dfrac{n}{\sqrt{4n^2+n}},\,\forall n\in{N^*}. 
		\end{eqnarray*}
		 Mà $\lim\dfrac{n}{\sqrt{4n^2}}=\lim\dfrac{1}{2}=\dfrac{1}{2};\,\lim\dfrac{n}{\sqrt{4n^2+n}}=\lim\dfrac{1}{\sqrt{4+\dfrac{1}{n}}}=\dfrac{1}{2}$.\\
	Nên $\lim\left(\dfrac{1}{\sqrt{4n^2+1}}+\dfrac{1}{\sqrt{4n^2+2}}+\cdots+\dfrac{1}{\sqrt{4n^2+n}}\right)=\dfrac{1}{2}$.} 
\end{ex}

% nguồn: Ok_TuanLuan_Toan11_Lan2
\begin{ex}%[1D3V1-3]%[ma:0de0d16c]
	 Tính giới hạn $\lim\left(\dfrac{1}{1\cdot3}+\dfrac{1}{3\cdot5}+\cdots+\dfrac{1}{\left(2n-1\right)\left(2n+1\right)}\right)$.
	\loigiai{ 
		 Do $\dfrac{1}{\left(2k-1\right)\left(2k+1\right)}=\dfrac{1}{2}\cdot\dfrac{2k+1-\left(2k-1\right)}{\left(2k-1\right)\left(2k+1\right)}=\dfrac{1}{2}\left(\dfrac{1}{2k-1}-\dfrac{1}{2k+1}\right)$.\\
		Suy  ra
		\begin{eqnarray*}
			&\Rightarrow&\lim\left(\dfrac{1}{1\cdot3}+\dfrac{1}{3\cdot5}+\cdots+\dfrac{1}{\left(2n-1\right)\left(2n+1\right)}\right)\\ 
			 &&=\lim\dfrac{1}{2}.\left(\dfrac{1}{1}-\dfrac{1}{3}+\dfrac{1}{3}-\dfrac{1}{5}+\cdots+\dfrac{1}{2n-1}-\dfrac{1}{2n+1}\right)=\lim\dfrac{1}{2}\left(1-\dfrac{1}{2n+1}\right)=\dfrac{1}{2}.
		\end{eqnarray*}
	} 
\end{ex}

% nguồn: Ok_TuanLuan_Toan11_Lan2
\begin{ex}%[1D3V1-3]%[ma:9888058e]
	Tìm  $\lim \left(2n-\sqrt{9n^2+n}+\sqrt{n^2+2n}\right)$.
	\loigiai{ 
		 \allowdisplaybreaks
		\begin{eqnarray*}
			\lim \left(2n-\sqrt{9n^2+n}+\sqrt{n^2+2n}\right) 
			 &=& \lim \left(3n-\sqrt{9n^2+n}+\sqrt{n^2+2n}-n\right).
		\end{eqnarray*}
		\begin{eqnarray*}
			\lim  \left(3n-\sqrt{9n^2+n}\right)
			&=& \lim \dfrac{\left(3n-\sqrt{9n^2+n}\right)\left(3n+\sqrt{9n^2+n}\right)}{3n+\sqrt{9n^2+n}} \\ 
			&=& \lim  \dfrac{-n}{3n+\sqrt{9n^2+n}} \\
			&=&  \lim \dfrac{-1}{3+\sqrt{9+\dfrac{1}{n}}} \\
			&=& -\dfrac{1}{6}. 
		\end{eqnarray*}
		\begin{eqnarray*}
			\lim \left(\sqrt{n^2+2n}-n\right) 
			&=& \lim \dfrac{\left(\sqrt{n^2+2n}-n\right)\left(\sqrt{n^2+2n}+n\right)}{\sqrt{n^2+2n}+n}  \\
			&=& \lim  \dfrac{2n}{\sqrt{n^2+2n}+n} \\
			&=& \lim  \dfrac{2}{\sqrt{1+\dfrac{2}{n}}+1} \\
			 &=& 1.
		\end{eqnarray*}
		Suy ra  $\lim \left(2n-\sqrt{9n^2+n}+\sqrt{n^2+2n}\right) = -\dfrac{1}{6}+1 = \dfrac{5}{6}$.
	 }
\end{ex}

% nguồn: Ok_TuanLuan_Toan11_Lan2
\begin{ex}%[1D3V1-3]%[ma:b73479d6]
	Tìm $\lim  \left(\sqrt[3]{8n^3+4n^2+2}-2n+3\right)$.
	\loigiai{ 
		 \allowdisplaybreaks
		\begin{eqnarray*}
			&&\lim  \left(\sqrt[3]{8n^3+4n^2+2}-2n+3\right)\\
			&=& \lim \left[\dfrac{\left(\sqrt[3]{8n^3+4n^2+2}-2n\right)\left({\sqrt[3]{\left(8n^3+4n^2+2\right)}}^2+2n\sqrt[3]{8n^3+4n^2+2}+4n^2\right)}{{\sqrt[3]{\left(8n^3+4n^2+2\right)}}^2+2n\sqrt[3]{8n^3+4n^2+2}+4n^2}+3\right]  \\
			&=& \lim \left[\dfrac{4n^2+2}{{\sqrt[3]{\left(8n^3+4n^2+2\right)}}^2+2n\sqrt[3]{8n^3+4n^2+2}+4n^2}+3\right] \\ 
			&=& \lim \left[\dfrac{4+\dfrac{2}{n^2}}{{\sqrt[3]{\left(8+\dfrac{4}{n}+\dfrac{2}{n^2}\right)}}^2+2\sqrt[3]{8+\dfrac{4}{n}+\dfrac{2}{n^2}}+4}+3\right] \\ 
			&=&  \dfrac{4}{12}+3 \\
			&=&  \dfrac{10}{3}.
		\end{eqnarray*}
	 }
\end{ex}

% nguồn: Ok_TuanLuan_Toan11_Lan2
\begin{ex}%[1D3V1-3]%[ma:7905ee95]
	 Tìm $\lim \left(\sqrt{9n^2+3n-4}-3n+2\right)$.
	 \loigiai{
		 \allowdisplaybreaks
		\begin{eqnarray*}
			\lim  \left(\sqrt{9n^2+3n-4}-3n+2\right)
			&=& \lim  \left[\dfrac{\left(\sqrt{9n^2+3n-4}-3n\right)\left(\sqrt{9n^2+3n-4}+3n\right)}{\sqrt{9n^2+3n-4}+3n}+2\right]\\
			&=& \lim \left[\dfrac{3n-4}{\sqrt{9n^2+3n-4}+3n}+2\right]\\ 
			&=& \lim  \left[\dfrac{3-\dfrac{4}{n}}{\sqrt{9+\dfrac{3}{n}-\dfrac{4}{n^2}}+3}+2\right] \\
			&=&  \dfrac{3}{6}+2\\
			&=& \dfrac{5}{2}. 
		\end{eqnarray*}
	} 
\end{ex}

% nguồn: Ok_TuanLuan_Toan11_Lan2
\begin{ex}%[1D3V1-3]%[ma:5d9b618e]
	Tính  giới hạn $\lim\dfrac{1\cdot3\cdot5\cdot7\cdots\left(2n-1\right)}{2\cdot4\cdot6\cdots\left(2n\right)}$
	\loigiai{ 
		\begin{itemize}
			\item \textbf{Cách  1.} Ta có
			\begin{eqnarray*}
				&u_n&=\dfrac{1\cdot3\cdot5\cdot7\cdots\left(2n-1\right)}{2\cdot4\cdot6\cdots\left(2n\right)}\\ 
				 &\Rightarrow u_n^2&=\dfrac{1^23^25^2\cdots\left(2n-1\right)^2}{2^24^26^2\cdots\left(2n\right)^2}\\
				 &&=\dfrac{1\cdot3}{2^2}\cdot\dfrac{3\cdot5}{4^2}\cdots\dfrac{\left(2n-1\right)\left(2n+1\right)}{\left(2n\right)^2}\cdot\dfrac{1}{2n+1}<\dfrac{1}{2n+1}.
			\end{eqnarray*}
			Vậy ta có $0 <u_n<\dfrac{1}{\sqrt{2n+1}}$,  $\forall n\in{N^*}$.\\
			 Mà $\lim\dfrac{1}{\sqrt{2n+1}}=0\Rightarrow\lim\dfrac{1\cdot3\cdot5\cdot7\cdots\left(2n-1\right)}{2\cdot 4\cdot6\cdots\left(2n\right)}=0$.\\
			\item \textbf{Cách 2.} Đặt $u_n=\dfrac{1}{2}\cdot\dfrac{3}{4}\cdot\dfrac{5}{6}\cdot\dfrac{7}{8}\cdots\dfrac{2n-1}{2n}$.\\ Ta có  $\dfrac{2k-1}{2k}=\dfrac{2k-1}{\sqrt{4k^2}}\le\dfrac{2k-1}{\sqrt{4k^2-1}}=\sqrt{\dfrac{2k-1}{2k+1}}$.\\
			Suy  ra $\heva{&\dfrac{1}{2}\le\sqrt{\dfrac{1}{3}}\\&\dfrac{3}{4}\le\sqrt{\dfrac{3}{5}}\\&\ldots\\&\dfrac{2n-1}{2n}\le\sqrt{\dfrac{2n-1}{2n+1}}}\Rightarrow\dfrac{1\cdot3\cdot5\cdot7\cdots\left(2n-1\right)}{2\cdot4\cdot6\cdot8\cdots2n}<\sqrt{\dfrac{1}{3}}\cdot\sqrt{\dfrac{3}{5}}\cdots\sqrt{\dfrac{2n-1}{2n+1}}=\dfrac{1}{\sqrt{2n+1}}$.\\
			 Suy ra $\left|u_n\right|\le\dfrac{1}{\sqrt{2n+1}}$ mà $\lim\dfrac{1}{\sqrt{2n+1}}=0$.
	\end{itemize}}
\end{ex}

% nguồn: Ok_TuanLuan_Toan11_Lan2
\begin{ex}%[1D3C1-3]%[ma:dd7682c4]
	Tìm  $\lim{u_n}$ biết $\left(u_n\right)\colon \heva{&{u_1}=2\\&{u_{n+1}}=\dfrac{u_n+1}{2}},\, n=1,2,3,\ldots$.
	 \loigiai{
	Tìm công thức  số hạng tổng quát là $u_n=\dfrac{2^{n-1}+1}{2^{n-1}}$ suy ra $\lim{u_n}=\lim\dfrac{2^{n-1}+1}{2^{n-1}}=1$.}
\end{ex}

% nguồn: Ok_TuanLuan_Toan11_Lan2
\begin{ex}%[1D3C1-3]%[ma:cbf02cd7]
	Tìm $\lim \left(\sqrt{n^4+n^2+1}-\sqrt[3]{n^6+1}\right)$. 
	\loigiai{ 
		Ta  có
		\allowdisplaybreaks 
		\begin{eqnarray*}
			 \lim \left(\sqrt{n^4+n^2+1}-\sqrt[3]{n^6+1}\right)
			 &=& \lim \left(\sqrt{n^4+n^2+1}-n^2+n^2-\sqrt[3]{n^6+1}\right).
		\end{eqnarray*}
		Ta có 
		\begin{eqnarray*}
			\lim \left(\sqrt{n^4+n^2+1}-n^2\right) 
			&=& \lim  \dfrac{\left(\sqrt{n^4+n^2+1}-n^2\right)\left(\sqrt{n^4+n^2+1}+n^2\right)}{\sqrt{n^4+n^2+1}+n^2} \\
			 &=& \lim \dfrac{n^2+1}{\sqrt{n^4+n^2+1}+n^2} \\
			&=& \lim  \dfrac{1+\dfrac{1}{n^2}}{\sqrt{1+\dfrac{1}{n^2}}+1} \\
			 &=& \dfrac{1}{2}.
		\end{eqnarray*}
		Ta  có
		\begin{eqnarray*}
			\lim  \left(n^2-\sqrt[3]{n^6+1}\right)
			&=& \lim  \dfrac{\left(n^2-\sqrt[3]{n^6+1}\right)\left(n^4+n^2\sqrt[3]{n^6+1}+\sqrt[3]{{\left(n^6+1\right)}^2}\right)}{n^4+n^2\sqrt[3]{n^6+1}+\sqrt[3]{{\left(n^6+1\right)}^2}} \\
			 &=& \lim \dfrac{-1}{n^4+n^2\sqrt[3]{n^6+1}+\sqrt[3]{{\left(n^6+1\right)}^2}} \\
			&=& \lim \dfrac{-\dfrac{1}{n^4}}{1+\sqrt[3]{1+\dfrac{1}{n^6}}+\sqrt[3]{{\left(1+\dfrac{1}{n^6}\right)}^2}} \\ 
			 &=& 0.
		\end{eqnarray*}
		Suy ra $\lim \left(\sqrt{n^4+n^2+1}-\sqrt[3]{n^6+1}\right) = \dfrac{1}{2} + 0 = \dfrac{1}{2}$. 
	 }
\end{ex}

% nguồn: Ok_TuanLuan_Toan11_Lan2
\begin{ex}%[1D3C1-3]%[ma:97324ab1]
	Tìm $\lim\dfrac{u_n}{n^2}$ biết  $\left(u_n\right)\colon \heva{&{u_1}=1,\, u_2=3\\
	&{u_{n+2}}=2u_{n+1}-u_n+1},\, n=1,2,3,\ldots$. 
	 \loigiai{
	Tìm công thức  số hạng tổng quát là $u_n=\dfrac{n\left(n+1\right)}{2}$ suy ra $\lim\dfrac{u_n}{n^2}=\dfrac{1}{2}$.}
\end{ex}

% nguồn: Ok_TuanLuan_Toan11_Lan2
\begin{ex}%[1D3C1-3]%[ma:2ddf3860]
	Tìm $\lim \left(\sqrt{n^2-2n}+2\sqrt[3]{n^2-8n^3}+3\sqrt{n^2+n}\right)$. 
	 \loigiai{
		 Ta có
		 \allowdisplaybreaks
		\begin{eqnarray*}
			 &&\lim \left(\sqrt{n^2-2n}+2\sqrt[3]{n^2-8n^3}+3\sqrt{n^2+n}\right)\\
			&=&  \lim \left[\left(\sqrt{n^2-2n}-n\right)+2\left(n+\sqrt[3]{n^2-8n^3}\right)+3\left(\sqrt{n^2+n}-n\right)\right].
		\end{eqnarray*}
		Ta  có
		\begin{eqnarray*}
			 \lim \left(\sqrt{n^2-2n}-n\right)
			 &=& \lim \dfrac{\left(\sqrt{n^2-2n}-n\right)\left(\sqrt{n^2-2n}+n\right)}{\sqrt{n^2-2n}+n} \\
			&=& \lim  \dfrac{-2n}{\sqrt{n^2-2n}+n} \\
			&=&  \lim \dfrac{-2}{\sqrt{1-\dfrac{2}{n}}+1} \\
			&=&  -1.
		\end{eqnarray*}
		Ta có 
		\begin{eqnarray*}
			&&\lim  \left(n+\sqrt[3]{n^2-8n^3}\right)\\
			&=& \lim \dfrac{\left(2n+\sqrt[3]{n^2-8n^3}\right)\left(4n^3-2n\sqrt[3]{n^2-8n^3}+\sqrt[3]{\left(n^2-8n^3\right)^2}\right)}{4n^2-2n\sqrt[3]{n^2-8n^3}+\sqrt[3]{\left(n^2-8n^3\right)^2}} \\ 
			&=& \lim \dfrac{n^2}{4n^3-2n\sqrt[3]{n^2-8n^3}+\sqrt[3]{\left(n^2-8n^3\right)^2}}  \\
			&=&  \lim \dfrac{1}{4-2\sqrt[3]{\dfrac{1}{n}-8}+\sqrt[3]{\left(\dfrac{1}{n}-8\right)^2}} \\
			&=&  \dfrac{1}{12}.
		\end{eqnarray*}
		Ta có 
		\begin{eqnarray*}
			\lim \left(\sqrt{n^2+n}-n\right) 
			&=& \lim  \dfrac{\left(\sqrt{n^2+n}-n\right)\left(\sqrt{n^2+n}+n\right)}{\sqrt{n^2+n}+n} \\
			&=& \lim \dfrac{n}{\sqrt{n^2+n}+n} \\ 
			&=&  \lim \dfrac{1}{\sqrt{1+\dfrac{1}{n}}+1} \\
			&=& \dfrac{1}{2}. 
		\end{eqnarray*}
		Suy ra $\lim \left(\sqrt{n^2-2n}+2\sqrt[3]{n^2-8n^3}+3\sqrt{n^2+n}\right) = -1 + 2 \cdot \dfrac{1}{12}  + 3 \cdot \dfrac{1}{2} = \dfrac{2}{3}$.
	 }
\end{ex}

% nguồn: Ok_TuanLuan_Toan11_Lan2
\begin{ex}%[1D3C1-3]%[ma:d2119f0c]
	Tìm $\lim{u_n}$  biết $\left(u_n\right)\colon \heva{&{u_1}=\dfrac{1}{2}\\&
	{u_{n+1}}=\dfrac{1}{2-u_n}},\, n=1,2,3,\ldots$. 
	\loigiai{ 
	Tìm công thức  số hạng tổng quát là $u_n=\dfrac{n}{n+1}$ suy ra $\lim{u_n}=\lim\dfrac{n}{n+1}=1$.}
\end{ex}

% nguồn: Ok_TuanLuan_Toan11_Lan2
\begin{ex}%[1D3C1-3]%[ma:0d47b9a7]
	Tìm $\lim  \left(\sqrt{4n^2+3n+7}-\sqrt[3]{8n^3+5n^2+1}\right)$.
	\loigiai{ 
		Ta có 
		\allowdisplaybreaks 
		\begin{eqnarray*}
			\lim  \left(\sqrt{4n^2+3n+7}-\sqrt[3]{8n^3+5n^2+1}\right)
			&=& \lim \left(\sqrt{4n^2+3n+7}-2n+2n-\sqrt[3]{8n^3+5n^2+1}\right). 
		\end{eqnarray*}
		Ta  có
		 \allowdisplaybreaks
		\begin{eqnarray*}
			\lim  \left(\sqrt{4n^2+3n+7}-2n\right)
			&=& \lim \dfrac{\left(\sqrt{4n^2+3n+7}-2n\right)\left(\sqrt{4n^2+3n+7}+2n\right)}{\sqrt{4n^2+3n+7}+2n}  \\
			&=& \lim \dfrac{3n+7}{\sqrt{4n^2+3n+7}+2n} \\ 
			&=& \lim \dfrac{3+\dfrac{7}{n}}{\sqrt{4+\dfrac{3}{n}+\dfrac{7}{n^2}}+2} \\ 
			&=& \dfrac{3}{4}. 
		\end{eqnarray*}
		Ta có 
		\begin{eqnarray*}
			&&\lim  \left(2n-\sqrt[3]{8n^3+5n^2+1}\right)\\
			&=& \lim  \dfrac{\left(2n-\sqrt[3]{8n^3+5n^2+1}\right)\left(4n^2+2n\sqrt[3]{8n^3+5n^2+1}+\sqrt[3]{{\left(8n^3+5n^2+1\right)}^2}\right)}{4n^2+2n\sqrt[3]{8n^3+5n^2+1}+\sqrt[3]{{\left(8n^3+5n^2+1\right)}^2}} \\
			 &=& \lim \dfrac{-5n^2-1}{4n^2+2n\sqrt[3]{8n^3+5n^2+1}+\sqrt[3]{{\left(8n^3+5n^2+1\right)}^2}} \\
			&=& \lim  \dfrac{-5-\dfrac{1}{n^2}}{4+2\sqrt[3]{8+\dfrac{5}{n}+\dfrac{1}{n^2}}+\sqrt[3]{{\left(8+\dfrac{5}{n}+\dfrac{1}{n^2}\right)}^2}} \\
			 &=& -\dfrac{5}{12}.
		\end{eqnarray*}
		Suy ra $\lim \left(\sqrt{4n^2+3n+7}-\sqrt[3]{8n^3+5n^2+1}\right) = \dfrac{3}{4} - \dfrac{5}{12}  = \dfrac{1}{3}$.
	} 
\end{ex}

% nguồn: Ok_TuanLuan_Toan11_Lan2
\begin{ex}%[1D3C1-3]%[ma:f4f604c3]
	Tìm $\lim\dfrac{u_n}{3\cdot 2^n}$ biết $\left(u_n\right)\colon \heva{&{u_1}=1,\, u_2=6\\ 
	&{u_{n+2}}-3u_{n+1}+2u_n=0},\,  n=1,2,3,\ldots$.
	\loigiai{ 
	Tìm công thức số hạng tổng quát là $u_n=-4+5\cdot 2^{n-1}$ suy ra $\lim\dfrac{u_n}{3\cdot  2^n}\lim \dfrac{-8+5\cdot 2^n}{6\cdot 2^n}=\dfrac{5}{6}$.}
\end{ex}
